% STAGE13-CHAPTER: V5-C04
% CHAPTER-SCHEMA-COVERAGE: CH-01,CH-02,CH-03,CH-04,CH-05,CH-06,CH-07,CH-08,CH-09,CH-10,CH-11,CH-12,CH-13,CH-14,CH-15,CH-16,CH-17,CH-18,CH-19,CH-20,CH-21,CH-22,CH-23,CH-24,CH-25,CH-26,CH-27,CH-28,CH-29,CH-30,CH-31,CH-32,CH-33,CH-34,CH-35,CH-36,CH-37
% CHAPTER-SCHEMA-STATUS: CH-01=passed; CH-02=passed; CH-03=passed; CH-04=passed; CH-05=passed; CH-06=passed; CH-07=passed; CH-08=passed; CH-09=passed; CH-10=passed; CH-11=passed; CH-12=passed; CH-13=passed; CH-14=passed; CH-15=passed; CH-16=passed; CH-17=passed; CH-18=passed; CH-19=passed; CH-20=passed; CH-21=passed; CH-22=passed; CH-23=passed; CH-24=passed; CH-25=passed; CH-26=passed; CH-27=passed; CH-28=passed; CH-29=passed; CH-30=passed; CH-31=passed; CH-32=passed; CH-33=passed; CH-34=passed; CH-35=passed; CH-36=passed; CH-37=passed
% CHAPTER-SCHEMA-EVIDENCE: CH-01=\label{struct:V5-C04-CH01}; CH-02=\label{struct:V5-C04-CH02}; CH-03=\label{struct:V5-C04-CH03}; CH-04=\label{struct:V5-C04-CH04}; CH-05=\label{struct:V5-C04-CH05}; CH-06=\label{struct:V5-C04-CH06}; CH-07=\label{struct:V5-C04-CH07}; CH-08=\label{struct:V5-C04-CH08}; CH-09=\label{struct:V5-C04-CH09}; CH-10=\label{struct:V5-C04-CH10}; CH-11=\label{struct:V5-C04-CH11}; CH-12=\label{struct:V5-C04-CH12}; CH-13=\label{struct:V5-C04-CH13}; CH-14=\label{struct:V5-C04-CH14}; CH-15=\label{struct:V5-C04-CH15}; CH-16=\label{struct:V5-C04-CH16}; CH-17=\label{struct:V5-C04-CH17}; CH-18=\label{struct:V5-C04-CH18}; CH-19=\label{struct:V5-C04-CH19}; CH-20=\label{struct:V5-C04-CH20}; CH-21=\label{struct:V5-C04-CH21}; CH-22=\label{struct:V5-C04-CH22}; CH-23=\label{struct:V5-C04-CH23}; CH-24=\label{struct:V5-C04-CH24}; CH-25=\label{struct:V5-C04-CH25}; CH-26=\label{struct:V5-C04-CH26}; CH-27=\label{struct:V5-C04-CH27}; CH-28=\label{struct:V5-C04-CH28}; CH-29=\label{struct:V5-C04-CH29}; CH-30=\label{struct:V5-C04-CH30}; CH-31=\label{struct:V5-C04-CH31}; CH-32=\label{struct:V5-C04-CH32}; CH-33=\label{struct:V5-C04-CH33}; CH-34=\label{struct:V5-C04-CH34}; CH-35=\label{struct:V5-C04-CH35}; CH-36=\label{struct:V5-C04-CH36}; CH-37=\label{struct:V5-C04-CH37}
% CHAPTER-MODULE-SEQUENCE: learning_navigation; strict_modeling; mathematical_development; multi_view_understanding; algorithm_engineering; examples_exercises; closure_self_check
% PROJECT_SPEC-V1.1-EXERCISE-LAYOUT: grouped; kept=12; source_group=1; supplement=11
% DERIVATION-CHECK: step-1; step-2; step-3
% FORMAL-RESULT-CHECK: results=4; proofs=4
% ALGORITHM-CHECK: algorithms=1; deterministic-contracts=1
% FIGURE-KNOWLEDGE-MAP: fig:V5-C04-dependency-graph -> SLM-19-05-003,SLM-19-05-004,PRE-SA-008,SLM-19-05-005,SLM-19-03-007,SLM-19-05-006
% FIGURE-KNOWLEDGE-MAP: fig:V5-C04-conditional-slice -> SLM-19-05-003,SLM-19-05-004
% FIGURE-KNOWLEDGE-MAP: fig:V5-C04-coordinate-sweep -> PRE-SA-008,SLM-19-05-005,SLM-19-03-007
% FIGURE-KNOWLEDGE-MAP: fig:V5-C04-gibbs-axis-path -> SLM-19-05-004,SLM-19-05-005
% FIGURE-KNOWLEDGE-MAP: fig:V5-C04-gibbs-vs-mh -> PRE-SA-008,SLM-19-05-005
% FIGURE-KNOWLEDGE-MAP: fig:V5-C04-bivariate-normal-conditionals -> SLM-19-05-004,SLM-19-03-007
% FIGURE-KNOWLEDGE-MAP: fig:V5-C04-bayes-markov-blanket -> SLM-19-05-006
% FIGURE-KNOWLEDGE-MAP: fig:V5-C04-mixing-rho-comparison -> PRE-SA-008,SLM-19-05-006

\chapter{Gibbs 抽样}\label{chap:V5-C04}
\index{Gibbs抽样}
\index{吉布斯抽样|see{Gibbs抽样}}

% KNOWLEDGE-PLAN: KN-V5-C33-PREREQUISITE-001 L1
\paragraph{读前自检：本章地图：Gibbs 抽样。}Gibbs 抽样章地图把“联合目标$\to$满条件概率核$\to$坐标更新$\to$扫描链与诊断”作为主线；请为其中首条箭头补出必要条件，并检查终点仍落在本章声明的输出域。\par

\SLChapterOpening{联合目标$\to$满条件概率核$\to$坐标更新$\to$扫描链与诊断}
{怎样从联合分布构造可测的满条件版本，并证明逐坐标更新保持目标？}
{读完可处理零边缘集合、证明坐标核可逆，并区分系统扫描的平稳性与可逆性。}
{联合/边缘/条件分布、可测核、MH细致平衡，以及有限链与Harris条件。}
{先复习\GlobalChapterRef{32}{5}{3}的MH零流与一般状态收敛层次，再进入满条件核。}
\phantomsection\label{struct:V5-C04-CH01}
% KNOWLEDGE-PLAN: KN-V5-C33-PREREQUISITE-002 L1
\paragraph{读前自检：本章可检验学习目标。}Gibbs 抽样的首个可检验目标是“从联合分布严格定义满条件分布，说明归一化分母、支持集和“正比于”记号的合法边界”；请从联合分布严格定义满条件分布，说明归一化分母、支持集和“正比于”记号的合法边界并写出中间结果，再按“中间结果逐项包含首目标“从联合分布严格定义满条件分布，说明归一化分母、支持集和“正比于”记号的合法边界”声明的对象、条件与结论”验收。\par

\chaptergoals{
\begin{itemize}[leftmargin=2em]
  \item 从联合分布严格定义满条件分布，说明归一化分母、支持集和“正比于”记号的合法边界；
  \item 把单坐标Gibbs更新写成马尔可夫核，证明它关于目标分布可逆并保持目标分布；
  \item 区分系统扫描与随机扫描，证明二者的平稳性，并说明系统扫描一般不再可逆；
  \item 从\MetropolisHastings{}接受比推导“精确满条件提议必接受”，同时处理零边缘质量与不精确条件抽样器；
  \item 手工执行二元正态Gibbs更新，推导相关系数对混合速度和有效样本量的影响；
  \item 按完整数学步骤实现逐变量抽样，检查中途停止、不可约性、数值异常与复杂度，不虚构随机运行结果。
\end{itemize}
}

\phantomsection\label{struct:V5-C04-CH02}
% KNOWLEDGE-PLAN: KN-V5-C33-PREREQUISITE-003 L1
\paragraph{读前自检：最短先修。}Gibbs 抽样的最短先修明确要求“本章直接使用所引章节的目标分布、马尔可夫核、细致平衡、\MetropolisHastings{}零流分支和样本诊断”；请把首项“本章直接使用所引章节的目标分布、马尔可夫核、细致平衡、\MetropolisHastings{}零流分支和样本诊断”中的第一个先修对象写成定义域--运算--输出三元组，并以“该三元组逐项满足首项“本章直接使用所引章节的目标分布、马尔可夫核、细致平衡、\MetropolisHastings{}零流分支和样本诊断”的对象类型与成立条件”判断能否进入本章。\par

\begin{prerequisitebox}
本章直接使用\GlobalChapterRef{32}{5}{3}的目标分布、马尔可夫核、细致平衡、\MetropolisHastings{}零流分支和样本诊断；还使用第1册概率基础中的联合、边缘、条件分布与条件独立。读者应能把联合密度积分成边缘密度，并能判断一个条件密度是否归一化。涉及连续变量时统一相对于乘积支配测度$\mu=\mu_1\otimes\cdots\otimes\mu_d$书写；有限状态情形把积分替换为求和。
\end{prerequisitebox}

\phantomsection\label{struct:V5-C04-CH03}
% KNOWLEDGE-PLAN: KN-V5-C33-PREREQUISITE-004 L1
\paragraph{读前自检：先修自检。}Gibbs 抽样先修自检固定核验“给定联合密度$\pi(x_1,x_2)$，能否计算边缘$\pi_2(x_2)$并写出$\pi(x_1\mid x_2)$”；请给定联合密度$\pi(x_1,x_2)$，实际计算边缘$\pi_2(x_2)$并写出$\pi(x_1\mid x_2)$并写出中间结果，再按“$\pi(x_1\mid x_2)$的计算或证明结果满足首项所列等式、定义域或判断”明确判为通过或失败。\par

\begin{precheckbox}
\begin{enumerate}[leftmargin=2.2em]
  \item 给定联合密度$\pi(x_1,x_2)$，能否计算边缘$\pi_2(x_2)$并写出$\pi(x_1\mid x_2)$？
  \item 能否解释条件密度只在分母为正的条件值上由比值唯一确定？
  \item 若两个核$K_1,K_2$分别保持$\pi$，能否逐步写出$\pi K_1K_2=\pi$？
  \item 能否区分“核保持$\pi$”“链不可约”和“时间平均收敛”三种结论？
  \item 能否把二次型$x_1^2-2\rho x_1x_2+x_2^2$对$x_1$配方？
\end{enumerate}
第1--2项不足时回看条件概率；第3--4项不足时回看第5册第1章；第5项是二元正态满条件分布的直接前置。
\end{precheckbox}

\phantomsection\label{struct:V5-C04-CH04}
% KNOWLEDGE-PLAN: KN-V5-C33-PREREQUISITE-005 L1
\paragraph{读前自检：依赖路线。}Gibbs 抽样依赖路线从“所引章节的MH可逆性、零流分支与一般状态收敛条件 $\to$ 本章的坐标Gibbs核”进入“联合目标$\pi(x_1,\ldots,x_d)$ $\to$ 每个坐标的满条件分布$\pi_j(\cdot\mid x_{-j})$ $\to$ 单坐标Gibbs核$K_j$ $\to$ 系统扫描乘积核或随机扫描混合核 $\to$ 平稳性与可达性检查 $\to$ 燃烧期后的相关样本与Monte Carlo诊断”；请写出箭头成立所需的定义域条件，并核对前者输出能作为后者输入。\par

\begin{dependencybox}
\GlobalChapterRef{32}{5}{3}的MH可逆性、零流分支与一般状态收敛条件 $\to$ 本章的坐标Gibbs核；
联合目标$\pi(x_1,\ldots,x_d)$
$\to$ 每个坐标的满条件分布$\pi_j(\cdot\mid x_{-j})$
$\to$ 单坐标Gibbs核$K_j$
$\to$ 系统扫描乘积核或随机扫描混合核
$\to$ 平稳性与可达性检查
$\to$ 燃烧期后的相关样本与Monte Carlo诊断；
概率图的局部分解
$\to$ Markov毯
$\to$ 只保留与当前变量相连的因子
$\to$ 可计算的局部满条件分布。
\end{dependencybox}
图\ref{fig:V5-C04-dependency-graph}概括从联合目标与局部因子到Monte Carlo诊断的六节点学习与计算主链。
% v2.7.0 figure UID: FIG-P630-01
% R2: six-node dependency chain; every segment is an independent directed edge.
\tikzset{slfig-FIG-P630-01/.style={font=\fontsize{9.6pt}{11.6pt}\selectfont,every path/.append style={line cap=round,line join=round}}}
\begin{figure}[htbp]
\centering
\begin{tikzpicture}[slfig-FIG-P630-01,>=Stealth,every node/.style={align=center},
  core/.style={draw=SLBlue,rounded corners=2.2pt,fill=SLBlue!6,
    minimum width=31mm,minimum height=13mm,inner xsep=1.5mm,inner ysep=1.2mm,
    font=\fontsize{9.6pt}{11.8pt}\selectfont,line width=.72pt},
  side/.style={draw=SLRuleGray,rounded corners=2.2pt,fill=SLSoftGray,
    text width=22.5mm,minimum height=16mm,inner xsep=1.5mm,inner ysep=1.0mm,
    font=\fontsize{9.6pt}{11.6pt}\selectfont,line width=.55pt},
  boundary/.style={draw=SLGold,rounded corners=2.2pt,fill=SLGold!8,
    minimum width=76mm,minimum height=9mm,inner xsep=2mm,inner ysep=1mm,
    font=\fontsize{10.0pt}{12.0pt}\selectfont\bfseries,line width=.72pt},
  flow/.style={draw=SLBlue,line width=.95pt,-{Stealth[length=2.0mm,width=1.35mm]}},
  leader/.style={draw=SLGray,line width=.50pt}]
  \node[core] (joint) at (-3.80,.95) {联合目标 / 局部因子};
  \node[core] (cond) at (0,.95) {给定 $x_{-j}$ 的满条件\\$\pi_j(\cdot\mid x_{-j})$};
  \node[core] (coord) at (3.80,.95) {单坐标核 $K_j$\\只更新 $x_j$};
  \node[core] (scan) at (3.80,-.95) {扫描核\\系统 / 随机};
  \node[core] (sample) at (0,-.95) {相关样本};
  \node[core] (diag) at (-3.80,-.95) {诊断\\MCSE / ESS / 轨迹};
  \draw[flow] (joint.east)--(cond.west);
  \draw[flow] (cond.east)--(coord.west);
  \draw[flow] (coord.south)--(scan.north);
  \draw[flow] (scan.west)--(sample.east);
  \draw[flow] (sample.west)--(diag.east);
  \node[side] (correct) at (-6.40,2.85) {正确性条件\\目标保持\\支持可达\\遍历性};
  \node[side] (mix) at (6.40,-2.70) {混合效率\\自相关长度\\有效样本量};
  \draw[leader] (correct.south east)--(joint.north west);
  \draw[leader] (mix.north west)--(scan.south east);
  \node[boundary] at (0,-3.10)
    {正确内核 $\ne$ 快速混合};
\end{tikzpicture}
\captionsetup{width=.94\linewidth}
\caption{满条件把联合目标转为单坐标更新，扫描后得到需以MCSE、ESS与轨迹诊断的相关样本}
\label{fig:V5-C04-dependency-graph}
\end{figure}

\FloatBarrier
\noindent\textbf{读图顺序。}先看给定其余坐标$x_{-j}$后，满条件$\pi_j(\cdot\mid x_{-j})$如何只重抽$x_j$；再看单坐标核$K_j$怎样组成系统或随机扫描核并产生相关样本；最终用MCSE、ESS与轨迹诊断评估统计误差和混合效率。主链箭头只表示学习与计算依赖，不是概率图的生成时间方向；目标保持仍不自动推出不可约、收敛或快速混合。\par

\phantomsection\label{struct:V5-C04-CH05}
% STAGE19-SYMBOL-INDEX-BEGIN: V5-C04
\index[symbols]{minus-mathcalx-minus-eq-minus-prod-minus-j-eq-1-d-minus-mathcalx-minus-j--s3c215bdb8e6a@$\mathcal X=\prod_{j=1}^d\mathcal X_j$：$d$个坐标或变量块的联合取值空间}
\index[symbols]{x-eq-x-1-minus-ldots-minus-x-d--s770e2781e4df@$x=(x_1,\ldots,x_d)$：一个完整链状态}
\index[symbols]{x-minus-j--s903ceb3775f8@$x_{-j}$：删除第$j$个坐标后的其余状态}
\index[symbols]{minus-pi-minus-x--s1e8c624b4f84@$\pi(x)$：Gibbs链的目标联合分布}
\index[symbols]{minus-pi-minus-minus-j-x-minus-j--sa33d569adf10@$\pi_{-j}(x_{-j})$：对$x_j$积分或求和得到的边缘}
\index[symbols]{minus-pi-minus-j-t-minus-midx-minus-minus-j--s94920e9bb480@$\pi_j(t\mid x_{-j})$：给定其余坐标时第$j$个坐标的满条件分布}
\index[symbols]{k-j--s1f845f3055a9@$K_j$：只重抽第$j$个坐标的Gibbs更新}
\index[symbols]{k-minus-mathrm-minus-sys--s04fed11466e2@$K_{\mathrm{sys}}$：一轮系统扫描$K_1K_2\cdots K_d$}
\index[symbols]{k-minus-mathrm-minus-rs--sef27a9d15e32@$K_{\mathrm{rs}}$：固定权重随机扫描$\sum_jw_jK_j$}
\index[symbols]{w-j--s4cf8c79f5476@$w_j$：随机扫描选择坐标$j$的概率}
\index[symbols]{m-n-t--sc166d5a6b8b1@$m,n,T$：燃烧轮数、保留轮数及$T=m+n$}
\index[symbols]{minus-rho-minus--sc248889be4d7@$\rho$：标准二元正态的相关系数}
\index[symbols]{minus-varepsilon-minus-jt--s255b930e6b90@$\varepsilon_{jt}$：二元正态Gibbs第$t$轮第$j$步的创新}
% STAGE19-SYMBOL-INDEX-END: V5-C04
\begin{symboltablebox}
\small
\begin{tabularx}{\linewidth}{@{}l l X@{}}
\toprule 符号 & 类型、维数或范围 & 含义\\ \midrule
$\mathcal X=\prod_{j=1}^d\mathcal X_j$ & 乘积状态空间 & $d$个坐标或变量块的联合取值空间\\
$x=(x_1,\ldots,x_d)$ & $\mathcal X$中元素 & 一个完整链状态\\
$x_{-j}$ & $\prod_{\ell\ne j}\mathcal X_\ell$中元素 & 删除第$j$个坐标后的其余状态\\
$\pi(x)$ & 概率密度或质量函数 & Gibbs链的目标联合分布\\
$\pi_{-j}(x_{-j})$ & 非负边缘密度 & 对$x_j$积分或求和得到的边缘\\
$\pi_j(t\mid x_{-j})$ & 条件密度 & 给定其余坐标时第$j$个坐标的满条件分布\\
$K_j$ & 马尔可夫核 & 只重抽第$j$个坐标的Gibbs更新\\
$K_{\mathrm{sys}}$ & 马尔可夫核 & 一轮系统扫描$K_1K_2\cdots K_d$\\
$K_{\mathrm{rs}}$ & 马尔可夫核 & 固定权重随机扫描$\sum_jw_jK_j$\\
$w_j$ & 正数且$\sum_jw_j=1$ & 随机扫描选择坐标$j$的概率\\
$m,n,T$ & 非负整数 & 燃烧轮数、保留轮数及$T=m+n$\\
$\rho$ & $(-1,1)$ & 标准二元正态的相关系数\\
$\varepsilon_{jt}$ & 标准正态随机量 & 二元正态Gibbs第$t$轮第$j$步的创新\\
\bottomrule
\end{tabularx}
\end{symboltablebox}

% KNOWLEDGE-PLAN: KN-V5-C33-FORMULA-001 L1
\paragraph{读前自检：Gibbs抽样的满条件分解。}Gibbs满条件由联合密度固定其余坐标后归一化得到；请对二元$\pi(x_1,x_2)$写出$\pi_1(x_1\mid x_2)$，核对系统扫描第二步必须使用刚更新的$x_1$而非旧值。\par

\SLReviewSection{完整条件分布}{sec:V5-C04-S01}
\knowledgeanchor{SLM-19-05-003}{Gibbs抽样的满条件分解}
\knowledgeanchor{SLM-19-05-004}{满条件分布的逐坐标更新原理}
\phantomsection\label{struct:V5-C04-CH06}
设目标联合密度$\pi(x_1,\ldots,x_d)$可以计算，但从整个$d$维目标直接抽样困难。若固定$x_{-j}$后的一维或低维条件分布容易抽样，就可以依次重抽各坐标，把高维问题拆成一系列局部问题。图\ref{fig:V5-C04-conditional-slice}具体采用零均值、单位方差、相关系数$\rho=3/5$的二元正态，并固定$a=1$、$b=4/5$；阅读任务是先看联合密度上的水平与竖直切片，再核对边缘分母$m_2(b)=\phi(4/5)$与$m_1(a)=\phi(1)$（$\phi$为标准正态密度），最后得到归一化的满条件密度。切片不是任意截面，必须在该条件值的边缘密度为正时归一化。
% v2.7.0 figure UID: FIG-P632-01
% One analytic model drives the joint contours, slice denominators and both
% normalized conditional densities.  No panel-level or whole-figure scaling.
\begingroup
\def\slpRhoNum{3}
\def\slpRhoDen{5}
\def\slpANum{1}
\def\slpADen{1}
\def\slpBNum{4}
\def\slpBDen{5}
\pgfmathsetmacro{\slpRho}{\slpRhoNum/\slpRhoDen}
\pgfmathsetmacro{\slpA}{\slpANum/\slpADen}
\pgfmathsetmacro{\slpB}{\slpBNum/\slpBDen}
\pgfmathsetmacro{\slpCondVar}{1-\slpRho*\slpRho}
\pgfmathsetmacro{\slpCondSd}{sqrt(\slpCondVar)}
\pgfmathsetmacro{\slpMeanOne}{\slpRho*\slpB}
\pgfmathsetmacro{\slpMeanTwo}{\slpRho*\slpA}
\pgfmathsetmacro{\slpCondPeak}{1/(sqrt(2*pi)*\slpCondSd)}
\pgfmathsetmacro{\slpMarginalTwo}{exp(-\slpB*\slpB/2)/sqrt(2*pi)}
\pgfmathsetmacro{\slpMarginalOne}{exp(-\slpA*\slpA/2)/sqrt(2*pi)}
\pgfmathsetmacro{\slpEigLong}{1+\slpRho}
\pgfmathsetmacro{\slpEigShort}{1-\slpRho}
\pgfmathtruncatemacro{\slpMeanOneNum}{\slpRhoNum*\slpBNum}
\pgfmathtruncatemacro{\slpMeanOneDen}{\slpRhoDen*\slpBDen}
\pgfmathtruncatemacro{\slpMeanTwoNum}{\slpRhoNum*\slpANum}
\pgfmathtruncatemacro{\slpMeanTwoDen}{\slpRhoDen*\slpADen}
\pgfmathtruncatemacro{\slpCondVarNum}{\slpRhoDen*\slpRhoDen-\slpRhoNum*\slpRhoNum}
\pgfmathtruncatemacro{\slpCondVarDen}{\slpRhoDen*\slpRhoDen}
\pgfmathtruncatemacro{\slpCondSdNum}{sqrt(\slpCondVarNum)}
\pgfmathtruncatemacro{\slpCondSdDen}{\slpRhoDen}
\pgfmathtruncatemacro{\slpCrossNum}{2*\slpRhoNum}
\pgfmathtruncatemacro{\slpJointNormDen}{2*\slpCondSdNum}
\pgfmathtruncatemacro{\slpJointQuadNum}{\slpRhoDen*\slpRhoDen}
\pgfmathtruncatemacro{\slpJointQuadDen}{2*\slpCondVarNum}
\def\slpCOuter{2.25}
\def\slpCMiddle{1.55}
\def\slpCInner{0.90}
\pgfmathsetmacro{\slpLongOuter}{\slpCOuter*sqrt(\slpEigLong)}
\pgfmathsetmacro{\slpShortOuter}{\slpCOuter*sqrt(\slpEigShort)}
\pgfmathsetmacro{\slpLongMiddle}{\slpCMiddle*sqrt(\slpEigLong)}
\pgfmathsetmacro{\slpShortMiddle}{\slpCMiddle*sqrt(\slpEigShort)}
\pgfmathsetmacro{\slpLongInner}{\slpCInner*sqrt(\slpEigLong)}
\pgfmathsetmacro{\slpShortInner}{\slpCInner*sqrt(\slpEigShort)}

\tikzset{
  slfig-FIG-P632-01/.style={every path/.append style={line cap=round,line join=round}},
  sl632-axis/.style={draw=SLTextGray,line width=0.74pt,->},
  sl632-contour-outer/.style={draw=SLRuleGray,line width=0.72pt,densely dotted},
  sl632-contour-mid/.style={draw=SLMidBlue,line width=0.88pt,dashed},
  sl632-contour-inner/.style={draw=SLDeepBlue,line width=1.02pt},
  sl632-slice-mid/.style={draw=SLMidBlue,line width=0.78pt,dash pattern=on 5pt off 2pt},
  sl632-slice-deep/.style={draw=SLDeepBlue,line width=0.82pt,dash pattern=on 1pt off 1.7pt},
  sl632-density-mid/.style={draw=SLMidBlue,line width=1.12pt},
  sl632-density-deep/.style={draw=SLDeepBlue,line width=1.12pt,dash pattern=on 5pt off 2pt},
  sl632-peak-mid/.style={draw=SLMidBlue,line width=0.58pt,dashed},
  sl632-peak-deep/.style={draw=SLDeepBlue,line width=0.58pt,densely dotted},
  sl632-map-mid/.style={draw=SLMidBlue,line width=0.92pt,dashed,->},
  sl632-map-deep/.style={draw=SLDeepBlue,line width=0.92pt,dash pattern=on 1pt off 1.7pt,->}
}
\pgfkeysifdefined{/tikz/alt}{}{\tikzset{alt/.code={}}}
\begin{figure}[htbp]
\centering
\begin{tikzpicture}[
  slfig-FIG-P632-01,
  >=Stealth,
  every node/.style={font=\fontsize{9.6pt}{11.5pt}\selectfont},
  >={Stealth[length=2.3mm,width=1.6mm]},
  alt={对象—关系—结论：同一零均值、单位方差、相关系数五分之三的二元正态在左侧形成主轴四十五度的椭圆等高线；水平虚线固定第二坐标为五分之四，竖直点线固定第一坐标为一。两条不交叉箭头分别表示原始联合截面除以边缘密度 m2 等于标准正态密度在五分之四处的值、m1 等于标准正态密度在一处的值，得到均值二十五分之十二与五分之三、方差二十五分之十六、全实线积分为一的两条满条件密度。零边缘分母只用预先指定的可测正则条件版本处理且仅边缘几乎处处唯一，本高斯例两个分母均为正。}]
  % Joint-density panel: every ellipse is c times the square roots of the
  % two covariance eigenvalues 1+rho and 1-rho; no aspect ratio is hand-set.
  \begin{scope}[shift={(0,0)}]
    \draw[sl632-axis] (-3.05,0) -- (3.08,0)
      node[below left,xshift=-1pt,yshift=-1pt] {$x_1$};
    \draw[sl632-axis] (0,-2.55) -- (0,2.15) node[above] {$x_2$};
    \begin{scope}[rotate=45]
      \draw[sl632-contour-outer] (0,0) ellipse
        [x radius=\slpLongOuter cm,y radius=\slpShortOuter cm];
      \draw[sl632-contour-mid] (0,0) ellipse
        [x radius=\slpLongMiddle cm,y radius=\slpShortMiddle cm];
      \draw[sl632-contour-inner] (0,0) ellipse
        [x radius=\slpLongInner cm,y radius=\slpShortInner cm];
    \end{scope}
    \draw[sl632-slice-mid] (-2.72,\slpB) -- (2.92,\slpB)
      node[pos=0.02,above=3pt] {$x_2=b=\dfrac{\slpBNum}{\slpBDen}$};
    \draw[sl632-slice-deep] (\slpA,-2.43) -- (\slpA,2.52);
    \draw[draw=SLDeepBlue,line width=0.58pt]
      (\slpA,2.52) -- (1.48,2.65);
    \node[anchor=west,fill=white,inner sep=5pt] at (1.58,2.72)
      {$x_1=a=\slpANum$};
    \fill[SLDeepBlue] (\slpA,\slpB) circle (2.2pt);
    \draw[draw=SLDeepBlue,line width=0.58pt]
      (\slpA,\slpB) -- (1.42,1.45);
    \node[fill=white,inner sep=5pt] at (1.72,1.70) {$(a,b)$};
    \node[align=center] at (0,3.92)
      {$\rho=\dfrac{\slpRhoNum}{\slpRhoDen},\quad a=\slpANum,\quad b=\dfrac{\slpBNum}{\slpBDen}$\\[-0.5pt]
       $\displaystyle \pi(x_1,x_2)=\frac{\slpRhoDen}{\slpJointNormDen\pi}
       \exp\!\left[-\frac{\slpJointQuadNum}{\slpJointQuadDen}q(x_1,x_2)\right]$\\[-0.5pt]
       $\displaystyle q=x_1^2-\frac{\slpCrossNum}{\slpRhoDen}x_1x_2+x_2^2$};
    \node[align=center,text=SLDeepBlue] at (0,-2.98)
      {联合等高线：主轴 $45^\circ$，半轴 $c\sqrt{1+\rho}$ 与 $c\sqrt{1-\rho}$};
  \end{scope}

  % X1 | X2=b panel.  The coordinate y unit is only an axis unit; text and
  % line widths are not scaled.
  \begin{scope}[shift={(7.08,1.20)}]
    \begin{scope}[y=3.1cm]
      \draw[sl632-axis] (-2.52,0) -- (2.78,0) node[right] {$t$};
      \draw[sl632-axis] (-2.42,-0.025) -- (-2.42,0.62);
      \draw[sl632-density-mid,domain=-2.34:2.66,samples=120,smooth]
        plot (\x,{\slpCondPeak*exp(-((\x-\slpMeanOne)^2)/(2*\slpCondVar))});
      \draw[sl632-peak-mid] (\slpMeanOne,0) -- (\slpMeanOne,\slpCondPeak);
      \node[below=2pt,text=SLMidBlue] at (\slpMeanOne,0)
        {$\dfrac{\slpMeanOneNum}{\slpMeanOneDen}$};
    \end{scope}
    \node[align=center] at (0.55,2.68)
      {$\displaystyle \pi_1(t\mid X_2=b)=\frac{\pi(t,b)}{m_2(b)}$\\[-0.5pt]
       $\displaystyle \mathcal N\!\left(\frac{\slpMeanOneNum}{\slpMeanOneDen},
       \frac{\slpCondVarNum}{\slpCondVarDen}\right),\quad
       m_2(b)=\phi\!\left(\frac{\slpBNum}{\slpBDen}\right)
       \approx\pgfmathprintnumber[fixed,precision=3]{\slpMarginalTwo}>0$\\[-0.5pt]
       $\displaystyle \int_{\mathbb R}\pi_1(t\mid b)\,\mathrm dt=1,\quad
       \max\pi_1=\frac{\slpCondSdDen}{\slpCondSdNum\sqrt{2\pi}}$};
  \end{scope}

  % X2 | X1=a panel, derived from the same variance and peak macros.
  \begin{scope}[shift={(7.08,-3.40)}]
    \begin{scope}[y=3.1cm]
      \draw[sl632-axis] (-2.52,0) -- (2.78,0) node[right] {$t$};
      \draw[sl632-axis] (-2.42,-0.025) -- (-2.42,0.62);
      \draw[sl632-density-deep,domain=-2.34:2.66,samples=120,smooth]
        plot (\x,{\slpCondPeak*exp(-((\x-\slpMeanTwo)^2)/(2*\slpCondVar))});
      \draw[sl632-peak-deep] (\slpMeanTwo,0) -- (\slpMeanTwo,\slpCondPeak);
      \node[below=2pt,text=SLDeepBlue] at (\slpMeanTwo,0)
        {$\dfrac{\slpMeanTwoNum}{\slpMeanTwoDen}$};
    \end{scope}
    \node[align=center] at (0.55,2.68)
      {$\displaystyle \pi_2(t\mid X_1=a)=\frac{\pi(a,t)}{m_1(a)}$\\[-0.5pt]
       $\displaystyle \mathcal N\!\left(\frac{\slpMeanTwoNum}{\slpMeanTwoDen},
       \frac{\slpCondVarNum}{\slpCondVarDen}\right),\quad
       m_1(a)=\phi(\slpANum)
       \approx\pgfmathprintnumber[fixed,precision=3]{\slpMarginalOne}>0$\\[-0.5pt]
       $\displaystyle \int_{\mathbb R}\pi_2(t\mid a)\,\mathrm dt=1,\quad
       \max\pi_2=\frac{\slpCondSdDen}{\slpCondSdNum\sqrt{2\pi}}$};
  \end{scope}

  \draw[sl632-map-mid] (2.92,\slpB) --
    node[pos=0.48,above=10pt,align=center]{水平截面\\[-1pt]$\div\,m_2(b)$}
    (3.72,1.20) -- (4.18,1.20);
  \draw[sl632-map-deep] (\slpA,-2.43) --
    node[pos=0.78,above=10pt,align=center]{竖直截面\\[-1pt]$\div\,m_1(a)$}
    (3.35,-2.43) -- (3.90,-3.40) -- (4.18,-3.40);

  \node[draw=SLRed,rounded corners=2pt,fill=SLRedBG,
    align=center,text width=6.5cm,inner sep=4pt] at (7.08,-5.42)
    {若边缘分母为 $0$：采用预先指定的可测正则条件版本，且仅在边缘几乎处处意义下唯一；本高斯例的两个分母均为正。};
\end{tikzpicture}
\caption{同一二元正态联合密度的两条截面除以相应边缘密度后，得到方差$16/25$、全实线积分为$1$的满条件密度；零边缘处须使用预先指定的正则条件版本。}
\label{fig:V5-C04-conditional-slice}
\end{figure}
\endgroup


\noindent\textbf{读图顺序。}先看左图固定$x_2=b=4/5$的水平切片与固定$x_1=a=1$的竖直切片；再看两条无交叉箭头如何分别用$m_2(b)=\phi(4/5)$与$m_1(a)=\phi(1)$归一化原始联合截面；得到$\mathcal N(12/25,16/25)$与$\mathcal N(3/5,16/25)$两条全实线积分为$1$的满条件密度。本高斯例两个分母均为正；一般情形分母为$0$时须采用预先指定的可测正则条件版本，且该版本仅在边缘几乎处处意义下唯一。\par
\FloatBarrier

\SLReadTranslation{“联合分布难以一次抽样”翻译成两个可执行问题：固定$x_2$后能否抽$x_1\mid x_2$，固定新$x_1$后能否抽$x_2\mid x_1$。Gibbs每次只解决其中一个问题。}
\SLMethodTrigger{看到联合目标的各个满条件分布容易抽样，就按“固定其余坐标、抽一个坐标、立即写回”的顺序构造一轮扫描。}

\begin{example}[先会执行：二元正态的一轮Gibbs]\label{exm:V5-C04-bridge-bvn}
设$(X_1,X_2)$为零均值、单位方差、相关系数$\rho\in(-1,1)$的二元正态，且
\[
X_1\mid X_2=x_2\sim\mathcal N(\rho x_2,1-\rho^2),
\qquad
X_2\mid X_1=x_1\sim\mathcal N(\rho x_1,1-\rho^2).
\]
从$x^{(0)}=(0,0)$出发，取固定标准正态创新$\varepsilon_1=1,\varepsilon_2=0$。执行一轮系统扫描Gibbs更新，求$x^{(1)}$，并说明第二步应使用旧值还是同轮新值。
\end{example}
\SLExampleSolutionHeading{exm:V5-C04-bridge-bvn}
\begin{solution}
\SLReadTranslation{题目要求完成“先会执行：二元正态的一轮Gibbs”：先写清目标分布、提议或满条件与支持，再构造转移/估计量并核验不变性或归一化。}
\SolGiven 目标分布、提议核或满条件在题面支持上有定义；所有条件分母和归一化常数为正，状态空间与随机变量取值域保持一致。
\SLMethodTrigger{先声明目标分布、提议/满条件与支持，再逐步构造转移或估计量并核验归一化。}
\SolDerive
\SolGiven 条件$\rho\in(-1,1)$保证$1-\rho^2>0$，两个满条件分布的标准差均为$\sqrt{1-\rho^2}$。

\SolPlan 系统扫描按坐标顺序抽样并立即写回，因此第二个坐标必须读取本轮刚得到的$x_1^{(1)}$。

\SolDerive 第一步使用旧的$x_2^{(0)}=0$：
\[
x_1^{(1)}=\sqrt{1-\rho^2},
\]
第二步立即使用该新值以及$\varepsilon_2=0$：
\[
x_2^{(1)}
=\rho x_1^{(1)}+\sqrt{1-\rho^2}\,\varepsilon_2
=\rho\sqrt{1-\rho^2}.
\]

\SolCheck 当$\rho=0$时结果为$(1,0)$，与两个独立标准正态坐标的条件更新一致；每一步的条件方差均为正。

\SolAnswer 一轮后的状态为
\[
x^{(1)}=\bigl(\sqrt{1-\rho^2},\,\rho\sqrt{1-\rho^2}\bigr).
\]
若第二步仍用$x_1^{(0)}$，得到的是并行更新，而不是本章定义的系统扫描。
\SolCheck 核对概率非负、归一化和条件分母为正；若是采样核，再检查目标分布不变性或支持连通性。
\end{solution}

\phantomsection\label{struct:V5-C04-CH07}
% KNOWLEDGE-PLAN: KN-V5-C33-DEFINITION-001 L1
\paragraph{读前自检：满条件分布。}满条件$\pi_j(t\mid x_{-j})$只在边缘归一化常数正且有限时由比值定义；请对$t$积分核对为1，并在零分母分支采用预先指定的条件核版本。\par

\begin{definition}[满条件分布]\label{def:V5-C04-full-conditional}
以下给出进阶严格版本。设$\mathcal X=\prod_{\ell=1}^d\mathcal X_\ell$，每个$(\mathcal X_\ell,\mathcal B_\ell)$都是非空标准Borel空间，$\mu_\ell$是$\sigma$有限测度，$\mu=\bigotimes_{\ell=1}^d\mu_\ell$，而$\pi$是关于$\mu$联合可测的概率密度。对固定$j$定义
\begin{equation}\label{eq:V5-C04-marginal-minus-j}
\pi_{-j}(x_{-j})
=\int_{\mathcal X_j}\pi(x_{-j},s)\,\mu_j(\mathrm ds).
\end{equation}
由Tonelli定理，$\pi_{-j}$可测；令$D_j=\{x_{-j}:0<\pi_{-j}(x_{-j})<\infty\}$。当$x_{-j}\in D_j$时，第$j$个坐标的满条件密度为
\begin{equation}\label{eq:V5-C04-full-conditional}
\pi_j(t\mid x_{-j})
=\frac{\pi(x_{-j},t)}{\pi_{-j}(x_{-j})},
\qquad t\in\mathcal X_j.
\end{equation}
它在$\mu_j$意义下积分为$1$。固定一个锚点$x_j^\circ\in\mathcal X_j$，并在所有条件值上定义概率核
\begin{equation}\label{eq:V5-C04-full-conditional-kernel}
\Pi_j(B\mid x_{-j})=
\begin{cases}
\displaystyle\frac{\int_B\pi(x_{-j},t)\,\mu_j(\mathrm dt)}{\pi_{-j}(x_{-j})},&x_{-j}\in D_j,\\[1.0em]
\delta_{x_j^\circ}(B),&x_{-j}\notin D_j,
\end{cases}
\qquad B\in\mathcal B_j.
\end{equation}
联合可测性与Tonelli定理保证$x_{-j}\mapsto\Pi_j(B\mid x_{-j})$可测，所以$\Pi_j$是正则条件分布的一个可测版本。集合$D_j^c$的$\pi_{-j}$质量为$0$；在该集合上换用任何概率核都不改变联合分解或目标保持性，且绝不能把式\eqref{eq:V5-C04-full-conditional}解释成$0/0$。
\end{definition}
\index{满条件分布}

\phantomsection\label{struct:V5-C04-CH08}
\textbf{模型组成与合法性。}一个可执行的Gibbs模型至少包含
\[
\bigl(\mathcal X,\mu,\pi,x^{(0)},\{\mathcal C_j\}_{j=1}^d,\mathcal R\bigr),
\]
其中$\pi$是可归一化的目标概率密度，$x^{(0)}$位于其支持内，$\mathcal C_j$表示从式\eqref{eq:V5-C04-full-conditional}精确抽样的步骤，$\mathcal R$提供彼此独立的随机数。只知道未归一化联合核$\widetilde\pi$并不妨碍推导满条件，因为与$t$无关的共同常数会在归一化时消去；但仍必须确认联合目标可归一化，并能对每个实际遇到的$x_{-j}$构造合法条件抽样器。

\phantomsection\label{struct:V5-C04-CH09}
\textbf{学习策略。}Gibbs抽样不是通过最小化损失得到一个点，而是构造保持整个目标分布的链。设计时依次检查：目标联合分布是否可归一化；满条件分布是否相容于同一个联合目标；单坐标移动是否能连接目标支持；更新顺序是否固定且明确；条件抽样是否精确；最后才比较分块方式、相关性和单轮成本。把条件众数代替条件随机抽样会得到坐标上升类算法，不再是Gibbs抽样。

\phantomsection\label{struct:V5-C04-CH10}
\textbf{学习算法。}在系统扫描中，第$t$轮从$x^{(t-1)}$开始，先抽$x_1^{(t)}$，再抽$x_2^{(t)}$；第$j$步使用本轮已经更新的$x_1^{(t)},\ldots,x_{j-1}^{(t)}$与尚未更新的$x_{j+1}^{(t-1)},\ldots,x_d^{(t-1)}$。图\ref{fig:V5-C04-coordinate-sweep}明确区分“轮内状态”和“轮末样本”。算法\ref{alg:V5-C04-systematic-gibbs}及其前置说明另行区分完整轮次与中途停止状态。
% v2.7.0 figure UID: FIG-P634-01
% Structured glyph-safe redraw: preserve the one-arrow/eight-slot semantics,
% but express low-profile punctuation and equality through grouping and arrows.
\tikzset{
  slfig-FIG-P634-01/.style={font=\fontsize{9.6pt}{11.5pt}\selectfont,every path/.append style={line cap=round,line join=round}},
  sl634-old/.style={draw=SLRuleGray,densely dotted,fill=white,minimum width=13.2mm,minimum height=9.5mm,line width=.78pt,inner sep=2pt},
  sl634-done/.style={draw=SLBlue,fill=SLBlue!4,pattern=north east lines,pattern color=SLRuleGray,minimum width=13.2mm,minimum height=9.5mm,line width=.95pt,inner sep=2pt},
  sl634-halo/.style={draw=none,fill=white,inner xsep=1.5pt,inner ysep=1.25pt},
  sl634-now/.style={draw=SLGold,fill=SLGold!9,minimum width=13.2mm,minimum height=9.5mm,line width=1.05pt,inner sep=2pt},
  sl634-card/.style={draw=SLRuleGray,rounded corners=1.8pt,fill=white,line width=.55pt,inner xsep=5pt}
}
\pgfkeysifdefined{/tikz/alt}{}{\tikzset{alt/.code={}}}
\begin{figure}[htbp]
\centering
\begin{tikzpicture}[slfig-FIG-P634-01,>=Stealth,every node/.style={font=\fontsize{9.6pt}{11.5pt}\selectfont,align=center},
  alt={对象—关系—结论：固定顺序为 1,2,…,j−1,j,j+1,…,d；当前第 j 步的 x^[j] 前段是本轮新值且后段是前轮旧值；x^[d]=x^(t) 仅在末位更新结束后成立，并记录为轮末样本}]
  \node[font=\fontsize{10.6pt}{12.7pt}\selectfont\bfseries,text=SLBlue] at (0,2.16) {单轮系统扫描坐标带};
  \node at (-5.2,1.48) {$1$};
  \node at (-3.7,1.48) {$2$};
  \node at (-2.2,1.48) {省略};
  \node at (-.7,1.48) {前位};
  \node at (.8,1.48) {当前};
  \node at (2.3,1.48) {后位};
  \node at (3.8,1.48) {省略};
  \node at (5.3,1.48) {末位};
  \draw[-{Stealth[length=1.8mm,width=1.25mm]},draw=SLGray,line width=.60pt] (-5.52,1.10)--(5.62,1.10)
    node[right=4pt,font=\fontsize{9.6pt}{11.5pt}\selectfont] {更新顺序};
  \node[sl634-done] at (-5.2,0) {};
  \node[sl634-done] at (-3.7,0) {};
  \node[sl634-done] at (-2.2,0) {};
  \node[sl634-done] at (-.7,0) {};
  \node[sl634-halo] at (-5.2,0) {坐标\\$1$};
  \node[sl634-halo] at (-3.7,0) {坐标\\$2$};
  \node[sl634-halo] at (-2.2,0) {中间\\坐标};
  \node[sl634-halo] at (-.7,0) {前段\\末位};
  \node[sl634-now] at (.8,0) {当前\\坐标};
  \node[sl634-old] at (2.3,0) {后段\\首位};
  \node[sl634-old] at (3.8,0) {中间\\坐标};
  \node[sl634-old] at (5.3,0) {末位\\坐标};
  \node[font=\fontsize{9.6pt}{11.5pt}\selectfont,text=SLBlue] at (-2.9,-.76) {本轮新值};
  \node[font=\fontsize{9.6pt}{11.5pt}\selectfont,text=SLGold] at (.8,-.76) {当前新值};
  \node[font=\fontsize{9.6pt}{11.5pt}\selectfont,text=SLTextGray] at (3.8,-.76) {前轮旧值};

  \node[sl634-card,minimum width=118mm,minimum height=10.5mm] at (0,-1.52) {};
  \node[font=\fontsize{10.0pt}{12.0pt}\selectfont] at (0,-1.34) {$x^{[j]}$ 当前子步状态};
  \node[font=\fontsize{9.8pt}{11.7pt}\selectfont,text=SLBlue] at (-2.65,-1.81)
    {起始至当前坐标\quad 本轮新值};
  \node[font=\fontsize{9.8pt}{11.7pt}\selectfont,text=SLTextGray] at (2.65,-1.81)
    {后续至末位坐标\quad 前轮旧值};

  \node[sl634-card,minimum width=123mm,minimum height=11mm] at (0,-2.72) {};
  \node[inner sep=2.5pt,font=\fontsize{10.0pt}{12.0pt}\selectfont] (sweepend) at (-2.65,-2.87) {$x^{[d]}$};
  \node[inner sep=2.5pt,font=\fontsize{10.0pt}{12.0pt}\selectfont] (roundstate) at (.85,-2.87) {$x^{(t)}$};
  \node[inner sep=2.5pt,font=\fontsize{9.8pt}{11.7pt}\selectfont,text=SLTextGray] (sample) at (4.25,-2.87) {轮末样本};
  \draw[<->,draw=SLGray,line width=.60pt] (sweepend.east)--(roundstate.west)
    node[midway,above=4pt,font=\fontsize{9.6pt}{11.5pt}\selectfont,text=SLTextGray] {状态相同};
  \draw[-{Stealth[length=1.5mm,width=1.05mm]},draw=SLGray,line width=.60pt] (roundstate.east)--(sample.west)
    node[midway,above=4pt,font=\fontsize{9.6pt}{11.5pt}\selectfont,text=SLTextGray] {仅此记录};
\end{tikzpicture}
\captionsetup{width=.94\linewidth}
\caption[系统扫描逐坐标即时写回并区分子步状态与轮末样本]{系统扫描按固定次序即时写回；当前子步的前段使用本轮新值，后段沿用前轮旧值；末位更新结束后，末位状态与本轮样本状态相同并记录为轮末样本。}
\label{fig:V5-C04-coordinate-sweep}
\end{figure}

\FloatBarrier
\noindent\textbf{读图顺序。}先按上方编号与箭头确认$j=1,\ldots,d$的固定次序；再看$x^{[j]}$左侧坐标已写入同轮新值、右侧坐标仍保留上一轮旧值，因而后续更新会读取此前的新值；最终只把$x^{[d]}=x^{(t)}$记作本轮样本，不把中间状态当作一轮输出。\par

% KNOWLEDGE-PLAN: KN-V5-C33-BOUNDARY-002 L1
\paragraph{读前自检：Gibbs抽样的满条件提议。}Gibbs坐标提议$K_j$从满条件更新$x_j$并保持$x_{-j}$不变；请写出$K_j(x,dy)=\Pi_j(dy_j\mid x_{-j})\delta_{x_{-j}}(dy_{-j})$，核对它是归一化核。\par
% KNOWLEDGE-PLAN: KN-V5-C33-ALGORITHM_IDEA-001 L1
\paragraph{读前自检：系统扫描Gibbs算法。}系统扫描Gibbs以合法初态、各坐标满条件抽样器和扫描预算为输入；请按$j=1,\ldots,d$依次用最新坐标状态更新一整轮，轮末原子提交，并在预算用尽时返回轨迹与诊断。\par
% KNOWLEDGE-PLAN: KN-V5-C33-ALGORITHM_IDEA-002 L1
\paragraph{读前自检：Gibbs坐标核的目标保持。}单坐标Gibbs核$K_j$只改变第$j$坐标；请把联合密度分解为边缘乘满条件，交换$x_j,y_j$核对双向测度不变，从而$\pi K_j=\pi$。\par

\SLTeachSection{逐变量Gibbs抽样}{sec:V5-C04-S02}
\knowledgeanchor{PRE-SA-008}{Gibbs抽样的满条件提议}
\knowledgeanchor{SLM-19-05-005}{系统扫描Gibbs算法}
\knowledgeanchor{SLM-19-03-007}{Gibbs坐标核的目标保持}
\phantomsection\label{struct:V5-C04-CH11}
% CORE-DERIVATION-PLAN: KN-V5-C33-DERIVATION-001
\SLKnowledgeBridge{用边缘--条件分解证明精确满条件更新保持联合目标分布。}{联合目标 $\pi$、除第 $j$ 坐标外的状态 $x_{-j}$ 与满条件核 $\Pi_j(x_{-j},dx_j)$。}{采用可测的满条件版本；零边缘质量集合上使用预先选定版本，不使用密度比。}{先写分解 $\pi(dx)=\pi_{-j}(dx_{-j})\Pi_j(x_{-j},dx_j)$。}

\begin{derivationgoalbox}
从联合目标$\pi$构造只更新坐标$j$的核，证明它保持目标分布，并说明精确满条件提议为什么无需\MetropolisHastings{}拒绝步骤。
\end{derivationgoalbox}

\begin{derivationprepbox}
使用定义\ref{def:V5-C04-full-conditional}构造的可测概率核$\Pi_j$、联合密度的边缘--条件分解和单坐标仅改变$x_j$这一事实。集合$D_j^c$具有零边缘质量，其上只使用预先选定的版本，不使用密度比。
\end{derivationprepbox}
由式\eqref{eq:V5-C04-full-conditional}，对$t$变化时与$t$无关的因子可以并入归一化常数，因此
\begin{equation}\label{eq:V5-C04-conditional-proportional}
\pi_j(t\mid x_{-j})\propto\pi(x_{-j},t),
\end{equation}
但右侧必须在$\mathcal X_j$上具有有限正积分。

\textbf{步骤1：定义局部核。}定义单坐标更新核
\begin{equation}\label{eq:V5-C04-coordinate-kernel}
K_j(x,\mathrm dy)
=\Pi_j(\mathrm dy_j\mid x_{-j})
\,\delta_{x_{-j}}(\mathrm dy_{-j}).
\end{equation}
因为$\Pi_j$是概率核，而$x\mapsto\delta_{x_{-j}}(B_{-j})$可测，式\eqref{eq:V5-C04-coordinate-kernel}先在可测矩形上满足核的可测性，再由单调类定理推广到乘积$\sigma$代数；故$K_j$确实是马尔可夫核。
\textbf{步骤2：分解联合目标。}若$y_{-j}=x_{-j}$，把联合密度分解为
\[
\pi(x)=\pi_{-j}(x_{-j})\pi_j(x_j\mid x_{-j}).
\]
\textbf{步骤3：比较双向流。}正向平稳流中与坐标$j$有关的部分为
\begin{equation}\label{eq:V5-C04-coordinate-flow}
\pi_{-j}(x_{-j})
\pi_j(x_j\mid x_{-j})
\pi_j(y_j\mid x_{-j}),
\end{equation}
交换$x_j,y_j$后式\eqref{eq:V5-C04-coordinate-flow}不变，所以单坐标核满足细致平衡。

\phantomsection\label{struct:V5-C04-CH12}
\begin{theorem}[单坐标Gibbs核的可逆性]\label{thm:V5-C04-coordinate-reversible}
在定义\ref{def:V5-C04-full-conditional}的标准Borel与$\sigma$有限条件下，用式\eqref{eq:V5-C04-full-conditional-kernel}的可测正则条件版本构造式\eqref{eq:V5-C04-coordinate-kernel}。则$K_j$关于$\pi$可逆，因而$\pi K_j=\pi$；在$D_j^c$上更换版本不改变该结论。
\end{theorem}

\phantomsection\label{struct:V5-C04-CH13}
% KNOWLEDGE-PLAN: KN-V5-C33-PROOF-001 L1
\paragraph{读前自检：证明：单坐标Gibbs核的可逆性。}单坐标Gibbs联合测度只在$x_{-j}=y_{-j}$上有质量，且含$\Pi_j(\mathrm dx_j\mid x_{-j})\Pi_j(\mathrm dy_j\mid x_{-j})$；请交换$x_j,y_j$核对测度不变，再积分一侧得到$\pi K_j=\pi$。\par

\begin{proof}
记$\mu_{-j}=\bigotimes_{\ell\ne j}\mu_\ell$。在$D_j$上有$\pi(\boldsymbol x)\mu(\mathrm d\boldsymbol x)=\pi_{-j}(x_{-j})\mu_{-j}(\mathrm dx_{-j})\Pi_j(\mathrm dx_j\mid x_{-j})$。在$D_j^c$上，外层边缘测度为$0$，所以任意版本贡献的联合质量都为$0$。因此在$\mathcal X\times\mathcal X$上，联合测度可写为
\[
\pi(\mathrm dx)K_j(x,\mathrm dy)
=\pi_{-j}(x_{-j})\mu_{-j}(\mathrm dx_{-j})
\Pi_j(\mathrm dx_j\mid x_{-j})
\Pi_j(\mathrm dy_j\mid x_{-j})
\delta_{x_{-j}}(\mathrm dy_{-j}).
\]
该测度只在$x_{-j}=y_{-j}$上有质量，并且两个条件概率因子的乘积在交换$x_j,y_j$后不变。因此
\[
\pi(\mathrm dx)K_j(x,\mathrm dy)
=\pi(\mathrm dy)K_j(y,\mathrm dx),
\]
即满足细致平衡。对$x$积分便得到$\pi K_j=\pi$。这也明确说明了零边缘版本为何不会进入任何正质量的双向流。
\end{proof}

\begin{proposition}[系统扫描与随机扫描]\label{prop:V5-C04-scan-invariance}
若每个$K_j$保持$\pi$，则系统扫描核
\[
K_{\mathrm{sys}}=K_1K_2\cdots K_d
\]
保持$\pi$。若$w_j>0$、$\sum_jw_j=1$且权重不依赖当前状态，则随机扫描核
\[
K_{\mathrm{rs}}=\sum_{j=1}^d w_jK_j
\]
关于$\pi$可逆。系统扫描核一般只保证平稳，不保证可逆。
\end{proposition}
% KNOWLEDGE-PLAN: KN-V5-C33-PROOF-002 L1
\paragraph{读前自检：证明：系统扫描与随机扫描。}系统扫描与随机扫描要求先确认“当这些核不交换时，正序与反序不同，不能由各分量可逆推出乘积可逆”；请随后把$\pi K_{\mathrm{sys}}$用到的关键定义展开并完成下一步变形，用乘以常数$w_j$再求和仍然对称，所以$K_{\mathrm{rs}}$可逆验收。\par

\begin{proof}
逐个使用$\pi K_j=\pi$得到
\[
\pi K_{\mathrm{sys}}
=(((\pi K_1)K_2)\cdots)K_d=\pi.
\]
对随机扫描，定理\ref{thm:V5-C04-coordinate-reversible}给出每个$K_j$的对称双向流；乘以常数$w_j$再求和仍然对称，所以$K_{\mathrm{rs}}$可逆。系统扫描的伴随核按反序排列为$K_d\cdots K_1$；当这些核不交换时，正序与反序不同，不能由各分量可逆推出乘积可逆。
\end{proof}

\begin{theorem}[有限全支持情形的遍历性]\label{thm:V5-C04-finite-positive}
设$\mathcal X=\prod_{j=1}^d\mathcal X_j$为有限乘积空间，且$\pi(x)>0$对每个$x\in\mathcal X$成立。则一轮系统扫描核$K_{\mathrm{sys}}$的每个矩阵元素都严格为正。因此该轮末链不可约、非周期，$\pi$是唯一平稳分布，并且从任意初值出发的边缘分布收敛到$\pi$。
\end{theorem}
% KNOWLEDGE-PLAN: KN-V5-C33-PROOF-003 L1
\paragraph{读前自检：证明：有限全支持情形的遍历性。}有限全支持Gibbs从$x$依次把前$j$个坐标改成$y$；请把一轮转移写成$\prod_{j=1}^d\pi_j(y_j\mid x_{-j}^{[j-1]})$，核对其严格为正，并令$y=x$验证不可约且周期为1。\par

\begin{proof}
任取$x,y\in\mathcal X$。令中间状态$x^{[j]}$的前$j$个坐标等于$y$，其余坐标等于$x$。从$x^{[j-1]}$更新第$j$个坐标到$y_j$的概率是相应满条件概率。全支持假设使联合质量与边缘质量都严格为正，所以每个满条件概率严格为正。于是
\[
K_{\mathrm{sys}}(x,y)
=\prod_{j=1}^d
\pi_j\!\left(y_j\mid x^{[j-1]}_{-j}\right)>0.
\]
故任意状态一步轮末转移即可到达任意状态，核不可约；取$y=x$又有$K_{\mathrm{sys}}(x,x)>0$，所以周期为$1$。有限状态马尔可夫链定理随后给出唯一平稳分布与边缘收敛，而命题\ref{prop:V5-C04-scan-invariance}已确认该平稳分布就是$\pi$。
\end{proof}

\begin{warningbox}
\textbf{定理边界与反例。}定理\ref{thm:V5-C04-finite-positive}的全支持条件很强。若目标支持受约束，单坐标移动可能无法连接不同部分。例如$\{(0,0),(1,1)\}$上的均匀分布具有合法但退化的满条件分布；从任一点更新单个坐标都不能离开原状态，Gibbs核等于恒等核。增加迭代轮数不能修复这种可约性。
\end{warningbox}

\SLTeachSection{与\MetropolisHastings{}的关系}{sec:V5-C04-S03}
\phantomsection\label{struct:V5-C04-CH14}
\textbf{几何解释。}在二维连续状态空间中，更新$x_1$时$x_2$固定，所以轨迹沿水平方向移动；随后更新$x_2$时$x_1$固定，轨迹沿竖直方向移动。图\ref{fig:V5-C04-gibbs-axis-path}把这种折线路径叠加在目标等密度线上。沿坐标轴移动不表示目标分布轴对齐；当等密度椭圆非常狭长且倾斜时，许多短的横纵步才能沿长轴前进。
% v2.6.0 figure UID: FIG-P637-01
% Semantic lock: exact contours, coordinates and six Gibbs moves are unchanged.
\tikzset{
  slfig-FIG-P637-01/.style={font=\fontsize{9.2pt}{11.0pt}\selectfont,every path/.append style={line cap=round,line join=round}},
  slfig-P637-axis/.style={draw=SLTextGray,line width=.72pt,->},
  slfig-P637-outer/.style={draw=SLRuleGray,line width=.72pt},
  slfig-P637-middle/.style={draw=SLMidBlue,line width=.80pt},
  slfig-P637-inner/.style={draw=SLDeepBlue,line width=.90pt},
  slfig-P637-xone/.style={->,draw=SLMidBlue,line width=1.00pt},
  slfig-P637-xtwo/.style={->,draw=SLDeepBlue,line width=1.00pt},
  slfig-P637-long/.style={-{Stealth[length=1.8mm]},draw=SLGold,line width=.78pt},
  slfig-P637-short/.style={-{Stealth[length=1.8mm]},draw=SLGray,line width=.70pt}
}
\pgfkeysifdefined{/tikz/alt}{}{\tikzset{alt/.code={}}}
\begin{figure}[!t]
\centering
\begin{tikzpicture}[slfig-FIG-P637-01,>=Stealth,every node/.style={font=\fontsize{9.2pt}{11.0pt}\selectfont},
  >={Stealth[length=2.3mm,width=1.6mm]},
  every node/.append style={inner sep=1.4pt},
  alt={对象—关系—结论：二维Gibbs轨迹在更新x下标1时水平移动、更新x下标2时竖直移动；折线是固定的教学示意而非伪造随机轨迹，倾斜狭长目标中的短轴向步揭示了强相关导致的慢混合}]
  \draw[slfig-P637-axis] (-4.3,0) -- (4.5,0) node[right=3pt] {$x_1$};
  \draw[slfig-P637-axis] (0,-3.4) -- (0,3.6) node[above=3pt] {$x_2$};

  \begin{scope}[rotate=34]
    \fill[SLSoftBlue,opacity=0.45] (0,0) ellipse (3.45 and 1.35);
    \draw[slfig-P637-outer] (0,0) ellipse (3.45 and 1.35);
    \draw[slfig-P637-middle] (0,0) ellipse (2.55 and 0.98);
    \draw[slfig-P637-inner] (0,0) ellipse (1.55 and 0.58);
  \end{scope}

  \coordinate (p0) at (-3.45,-1.8);
  \coordinate (p1) at (-1.15,-1.8);
  \coordinate (p2) at (-1.15,-0.20);
  \coordinate (p3) at (0.85,-0.20);
  \coordinate (p4) at (0.85,1.15);
  \coordinate (p5) at (2.35,1.15);
  \coordinate (p6) at (2.35,2.05);

  \draw[slfig-P637-xone] (p0) -- node[pos=.24,below=20pt,text=SLMidBlue] {更新$x_1$} (p1);
  \draw[slfig-P637-xtwo] (p1) -- node[pos=.02,below=24pt,text=SLDeepBlue] {更新$x_2$} (p2);
  \draw[slfig-P637-xone] (p2) -- (p3);
  \draw[slfig-P637-xtwo] (p3) -- (p4);
  \draw[slfig-P637-xone] (p4) -- (p5);
  \draw[slfig-P637-xtwo] (p5) -- (p6);

  \foreach \p/\lab/\anchor/\xs/\ys in {p0/$0$/north east/-2pt/-3pt,p1/$1$/south west/5pt/12pt,p2/$2$/north east/-10pt/-3pt,p3/$3$/north west/14pt/-20pt,p4/$4$/north west/17pt/-11pt,p5/$5$/north west/3pt/-3pt,p6/$6$/east/-7pt/-5pt}{
    \fill[SLDeepBlue] (\p) circle (2.2pt) node[anchor=\anchor,xshift=\xs,yshift=\ys,inner sep=0pt,font=\fontsize{8.8pt}{10.2pt}\selectfont] {\lab};
  }
  \draw[slfig-P637-long] (-2.2,-1.25)--(2.15,1.70)
    node[pos=.30,above=42pt,sloped,text=SLGold] {长轴};
  \draw[slfig-P637-short] (-.55,.80)--(.55,-.80)
    node[pos=.86,below right=40pt,text=SLGray] {短轴};
  \node[draw=SLRuleGray,fill=white,rounded corners=2pt,align=center,
    line width=.65pt,inner xsep=5pt,inner ysep=4pt,text width=52mm] at (0,-4.05)
    {固定教学示意：每次只改变一个坐标；轴向短步使沿长轴推进缓慢。};
\end{tikzpicture}
\caption{二维Gibbs轨迹在更新$x_1$时水平移动、更新$x_2$时竖直移动；折线是固定的教学示意而非伪造随机轨迹，倾斜狭长目标中的短轴向步揭示了强相关导致的慢混合}
\label{fig:V5-C04-gibbs-axis-path}
\end{figure}


\phantomsection\label{struct:V5-C04-CH15}
\textbf{概率解释。}更新坐标$j$时，算法保留$x_{-j}$，并用其正确条件分布中的新抽样替换旧$x_j$。给定$x_{-j}$后，新旧坐标是同一条件分布的两个取值，式\eqref{eq:V5-C04-coordinate-flow}因此天然对称。若把满条件分布作为单坐标\mbox{\MetropolisHastings{}}提议，则在正流分支上
\begin{align}
R(x,y)
&=\frac{\pi(y)\pi_j(x_j\mid y_{-j})}
        {\pi(x)\pi_j(y_j\mid x_{-j})}\\
&=\frac{\pi_{-j}(x_{-j})\pi_j(y_j\mid x_{-j})
          \pi_j(x_j\mid x_{-j})}
        {\pi_{-j}(x_{-j})\pi_j(x_j\mid x_{-j})
          \pi_j(y_j\mid x_{-j})}=1,
\end{align}
因为$y_{-j}=x_{-j}$。所以精确满条件提议必接受。图\ref{fig:V5-C04-gibbs-vs-mh}同时说明：若只近似满条件或改用其他提议，接受率不再自动等于$1$，此时必须恢复\MetropolisHastings{}校正。
% v2.6.0 figure UID: FIG-P638-01
% Semantic lock: four nodes, four draw paths, formulas and directions unchanged.
\tikzset{
  slfig-FIG-P638-01/.style={font=\fontsize{9.2pt}{11.0pt}\selectfont,every path/.append style={line cap=round,line join=round}},
  slfig-P638-eq/.style={minimum width=36mm,minimum height=13mm,inner xsep=1.5mm,inner ysep=1.6mm,outer sep=1.4pt,fill=white},
  slfig-P638-flow/.style={-{Stealth[length=2.0mm]},draw=SLBlue,line width=.95pt},
  slfig-P638-separator/.style={draw=SLRuleGray,line width=.50pt},
  slfig-P638-warn/.style={-{Stealth[length=1.8mm]},draw=SLRed,line width=.72pt,densely dashed},
  slfig-P638-exception/.style={draw=SLRed,rounded corners=2pt,fill=SLRedBG,line width=.72pt,minimum width=78mm,minimum height=13mm,inner xsep=5pt,inner ysep=4pt,outer sep=.8pt}
}
\begin{figure}[htbp]
\centering
\begin{tikzpicture}[slfig-FIG-P638-01,>=Stealth,every node/.style={font=\fontsize{9.2pt}{11.0pt}\selectfont,align=center}]
  \node[slfig-P638-eq] (q) at (-4.2,0) {{\color{SLBlue}\bfseries 1}\quad 精确满条件提议\\$q_j=\pi(x_j\mid x_{-j})$};
  \node[slfig-P638-eq] (r) at (0,0) {{\color{SLBlue}\bfseries 2}\quad MH 比值逐项抵消\\$R=\dfrac{\pi(y)\pi_j(x_j\mid x_{-j})}{\pi(x)\pi_j(y_j\mid x_{-j})}=1$};
  \node[slfig-P638-eq] (a) at (4.2,0) {{\color{SLBlue}\bfseries 3}\quad $\alpha=1$\\直接接受 $x_j\leftarrow y_j$};
  \draw[slfig-P638-flow] (q)--(r)--(a);
  \draw[slfig-P638-separator] (-6.1,-.85)--(6.1,-.85);
  \node[slfig-P638-exception] (w) at (0,-2.0)
    {近似满条件 / 其他提议 $q_j$\\恢复 MH 接受率校正；拒绝时保持 $x_j$（自环）};
  \draw[slfig-P638-warn] (q.south)--(w.north west);
  \draw[slfig-P638-warn] (a.south)--(w.north east);
\end{tikzpicture}
\captionsetup{width=.94\linewidth}
\caption{精确满条件分布作为单坐标提议时，正流分支的Metropolis--Hastings比值恒为1，Gibbs更新无需拒绝；若提议只近似满条件或改用其他分布，则必须恢复接受率校正和拒绝自环}
\label{fig:V5-C04-gibbs-vs-mh}
\end{figure}


\phantomsection\label{struct:V5-C04-CH16}
\textbf{优化解释及不适用边界。}标准Gibbs抽样没有要逐轮单调下降的目标函数，因此“优化收敛率”“局部最优点”和“每步改进损失”不适用于本算法。若把条件抽样改成条件众数，会得到迭代条件众数或坐标上升方法；其输出是点估计，不能替代后验样本。这里唯一可作设计优化的是分块、参数化与实现成本，而不是把随机状态更新解释成确定性下降。

\phantomsection\label{struct:V5-C04-CH17}
\textbf{图形辅助。}图\ref{fig:V5-C04-coordinate-sweep}强调轮内新旧坐标的混合使用；图\ref{fig:V5-C04-gibbs-axis-path}强调坐标方向路径；图\ref{fig:V5-C04-gibbs-vs-mh}强调精确条件提议与一般提议的接受规则。三图分别服务于更新顺序、几何混合与算法边界，不能互相替代。

\phantomsection\label{struct:V5-C04-CH18}
\begin{proposition}[二元正态Gibbs的边缘递推]\label{prop:V5-C04-bivariate-normal}
设
\[
(X_1,X_2)^\mathsf T\sim
N\!\left(0,
\begin{pmatrix}1&\rho\\ \rho&1\end{pmatrix}\right),
\qquad |\rho|<1.
\]
按$x_1$后$x_2$的顺序系统扫描，并令$\sigma=\sqrt{1-\rho^2}$。则
\begin{align}
X_1^{(t)}&=\rho X_2^{(t-1)}+\sigma\varepsilon_{1t},\\
X_2^{(t)}&=\rho X_1^{(t)}+\sigma\varepsilon_{2t}
=\rho^2X_2^{(t-1)}+\xi_t,
\end{align}
其中$\varepsilon_{1t},\varepsilon_{2t}$独立同分布于$N(0,1)$，$\xi_t$独立于过去且$\operatorname{Var}(\xi_t)=1-\rho^4$。在平稳状态下，$X_2$的滞后$k$自相关为$\rho^{2k}$，从而
\[
\tau_{\mathrm{int}}=\frac{1+\rho^2}{1-\rho^2},
\qquad
\frac{N_{\mathrm{eff}}}{N}=\frac{1-\rho^2}{1+\rho^2}.
\]
\end{proposition}
% KNOWLEDGE-PLAN: KN-V5-C33-PROOF-004 L1
\paragraph{读前自检：证明：二元正态Gibbs的边缘递推。}二元正态Gibbs按$x_1$后$x_2$扫描且$|\rho|<1$；请把两条满条件更新代入，得到$X_2^{(t)}=\rho^2X_2^{(t-1)}+\xi_t$，再算$\operatorname{Var}(\xi_t)=1-\rho^4$并核对滞后相关为$\rho^{2k}$。\par

\begin{proof}
联合密度指数中的二次型满足
\[
x_1^2-2\rho x_1x_2+x_2^2
=(x_1-\rho x_2)^2+(1-\rho^2)x_2^2.
\]
固定$x_2$后，与$x_1$有关的部分给出
$X_1\mid X_2=x_2\sim N(\rho x_2,1-\rho^2)$；交换坐标得到
$X_2\mid X_1=x_1\sim N(\rho x_1,1-\rho^2)$。代入两步更新便得到边缘递推，其中
\[
\xi_t=\sigma(\rho\varepsilon_{1t}+\varepsilon_{2t}),
\qquad
\operatorname{Var}(\xi_t)
=(1-\rho^2)(1+\rho^2)=1-\rho^4.
\]
这是系数为$\rho^2$的平稳AR(1)递推，所以$\operatorname{Corr}(X_2^{(t)},X_2^{(t+k)})=\rho^{2k}$。对几何级数求和得到
$1+2\sum_{k\ge1}\rho^{2k}=(1+\rho^2)/(1-\rho^2)$，再按ESS定义取倒数比例。
\end{proof}

图\ref{fig:V5-C04-bivariate-normal-conditionals}显示两个正态满条件切片。
% v2.3.1 figure UID: FIG-P639-01
% v2.6.0 redraw: preserve both exact densities while separating every annotation from data.
\tikzset{
  slfig-FIG-P639-01/.style={every path/.append style={line cap=round,line join=round}},
  slfig-FIG-P639-01-label/.style={font=\fontsize{9.2}{11}\selectfont},
  slfig-FIG-P639-01-note/.style={draw=SLRuleGray,line width=.55pt,rounded corners=2pt,
    fill=white,inner xsep=6pt,inner ysep=4pt,outer sep=1.2pt,align=center,
    font=\fontsize{9.2}{11}\selectfont}
}
\pgfplotsset{slfig-FIG-P639-01-axis/.style={
  tick label style={font=\fontsize{8.5}{10}\selectfont},
  label style={font=\fontsize{9.2}{11}\selectfont},clip=false}}
\begin{figure}[htbp]
\centering
\begin{tikzpicture}[slfig-FIG-P639-01]
\begin{axis}[slfig-FIG-P639-01-axis,width=10.5cm,height=5.8cm,xmin=-2,xmax=3,ymin=0,ymax=.56,
  axis lines=left,xlabel={$t$},ylabel={密度},xtick={-2,-1,0,1,2,3},ytick={0,.25,.5},
  axis line style={line width=.65pt},tick style={line width=.55pt},clip=false]
  \addplot[draw=SLBlue,fill=SLBlue!7,line width=1pt,domain=-2:3,samples=321]
    {0.498678*exp(-0.78125*(x-.45)^2)} \closedcycle;
  \addplot[draw=SLGold,line width=.95pt,densely dashed,domain=-2:3,samples=321]
    {0.498678*exp(-0.78125*(x-.60)^2)};
  \draw[draw=SLBlue,line width=.55pt,densely dashed] (axis cs:.45,0)--(axis cs:.45,.499);
  \draw[draw=SLGold,line width=.55pt,dash dot] (axis cs:.60,0)--(axis cs:.60,.499);
  \node[slfig-FIG-P639-01-label,anchor=south east,text=SLBlue] at (axis cs:.27,.556)
    {$N(0.45,0.64)$\quad $\mu_1=.45$};
  \node[slfig-FIG-P639-01-label,anchor=south west,text=SLGold] at (axis cs:.78,.556)
    {$N(0.60,0.64)$\quad $\mu_2=.60$};
  \node[slfig-FIG-P639-01-note,anchor=west] at (axis cs:3.14,.28)
    {共同方差 $0.64$\\均值随另一坐标改变};
\end{axis}
\end{tikzpicture}
\caption{取$\rho=0.6$、$a=1$、$b=0.75$时，两个满条件分布分别为$N(0.45,0.64)$与$N(0.60,0.64)$；图中曲线按同一组参数直接计算。}
\label{fig:V5-C04-bivariate-normal-conditionals}
\end{figure}
\FloatBarrier

% v2.6 layout: former forced clear removed; natural pagination.
图\ref{fig:V5-C04-mixing-rho-comparison}使用上述解析自相关而非伪造轨迹，比较不同$|\rho|$下的混合速度。
% v2.7.0 figure UID: FIG-P640-01
% v2.7.0 redraw: separate ACF from asymptotic ESS and preserve the |rho|<1 boundary.
\tikzset{slfig-FIG-P640-01/.style={every path/.append style={line cap=round,line join=round}}}
\pgfplotsset{slfig-FIG-P640-01-axis/.style={
  tick label style={font=\fontsize{9.6}{11.5}\selectfont},
  label style={font=\fontsize{9.8}{12}\selectfont},
  title style={font=\fontsize{9.6}{11.5}\selectfont,align=center},
  axis line style={draw=SLTextGray,line width=.65pt},
  tick style={draw=SLTextGray,line width=.55pt},clip=false}}
\begin{figure}[htbp]
\centering
\begin{tikzpicture}[slfig-FIG-P640-01]
\begin{axis}[slfig-FIG-P640-01-axis,
  name=main,width=8.6cm,height=5.7cm,
  xmin=0,xmax=12,
  ymin=0,ymax=1.05,
  xtick={0,2,4,6,8,10,12},
  ytick={0,0.25,0.5,0.75,1},
  title={(a) 轮末 ACF：$\rho^{2k}$},
  xlabel={滞后 $k$},
  ylabel={$\operatorname{Corr}(X_2^{(t)},X_2^{(t+k)})=\rho^{2k}$},
  legend style={draw=none,fill=none,font=\fontsize{9.6}{11.5}\selectfont,
    at={(0.5,-.28)},anchor=north,legend columns=3,/tikz/every even column/.append style={column sep=7pt}},
  mark size=1.7pt,samples=181,
  domain=0:12
]
\addplot[draw=SLBlue,line width=1pt] {pow(0.9025,x)};
\addlegendentry{$|\rho|=.95$}
\addplot[draw=SLTeal,line width=.85pt,densely dashed] {pow(0.49,x)};
\addlegendentry{$|\rho|=.70$}
\addplot[draw=SLGray,line width=.8pt,dash dot] {pow(0.04,x)};
\addlegendentry{$|\rho|=.20$}
\end{axis}
\begin{axis}[slfig-FIG-P640-01-axis,
  at={(main.outer north east)},anchor=outer north west,xshift=9mm,
  width=4.9cm,height=4.2cm,xmin=0,xmax=.99,ymin=-.06,ymax=1,
  title={(b) 渐近 ESS 比例：\\[-1pt]$\frac{1-\rho^2}{1+\rho^2}$},
  xlabel={$|\rho|$},ylabel={$N_{\rm eff}/N$},
  xtick={0,.5,.99},xticklabels={0,.5,.99},ytick={0,.5,1},
  axis lines=left,clip=false]
  \addplot[draw=SLGold,line width=1pt,domain=0:.99,samples=181] {(1-x^2)/(1+x^2)};
  \addplot[only marks,mark=*,mark size=1.8pt,draw=SLGold,fill=white]
    coordinates {(.99,0.0100499975)};
  \node[anchor=south east,xshift=-2pt,yshift=2pt,
    font=\fontsize{9.6}{11.5}\selectfont,text=SLGold,fill=white,inner sep=1pt]
    at (axis cs:.99,0.0100499975) {$(.99,\,.010)$};
  \node[anchor=north east,font=\fontsize{9.6}{11.5}\selectfont,text=SLGold,align=right,fill=white,inner sep=1pt]
    at (axis description cs:.98,.96) {$|\rho|\to1^-:$\\$N_{\rm eff}/N\to0$};
\end{axis}
\end{tikzpicture}
\captionsetup{width=.94\linewidth}
\caption{二元正态系统Gibbs的轮末解析自相关为$\rho^{2k}$，渐近有效样本比例为$(1-\rho^2)/(1+\rho^2)$；当$|\rho|$接近1时相关性衰减缓慢，即使内核正确且接受率为1，固定长度链的统计效率仍会显著下降}
\label{fig:V5-C04-mixing-rho-comparison}
\end{figure}

\FloatBarrier
读图时应先看左面板中轮末自相关$\rho^{2k}$随滞后$k$的衰减，再看右面板中渐近有效样本比例$(1-\rho^2)/(1+\rho^2)$随$|\rho|$的下降；当$|\rho|\to1^-$时，前者衰减趋于停滞而后者趋于$0$，这只是从合法区域$|\rho|<1$内逼近边界的极限，并不表示可以取退化相关系数$|\rho|=1$。

\Needspace{6\baselineskip}
\begin{example}[两轮确定性更新测试]\label{exm:V5-C04-two-sweeps}
取$\rho=1/2$、$\sigma=\sqrt3/2$、初值$x^{(0)}=(2,-1)$。把两轮标准正态输入固定为
\[
(\varepsilon_{11},\varepsilon_{21})=(0,1),
\qquad
(\varepsilon_{12},\varepsilon_{22})=(-1,0).
\]
按“先$X_1$后$X_2$”的系统扫描顺序复算两轮轮末状态，并说明该确定性测试能验证什么。
\end{example}

\Needspace{4\baselineskip}
\SLExampleSolutionHeading{exm:V5-C04-two-sweeps}
\begin{solution}
\SLReadTranslation{给定创新量把两轮随机更新固定成可逐步复算的确定性测试；每轮必须先得到新$x_1^{(t)}$，再用它计算同轮$x_2^{(t)}$。}
\SolGiven $\rho=1/2$、$\sigma=\sqrt3/2$、初值和四个创新量均已固定；本题只核验顺序更新公式，不用两轮数值判断随机链收敛。
\SLMethodTrigger{逐轮执行$x_1^{(t)}=\rho x_2^{(t-1)}+\sigma\varepsilon_{1t}$，立即把新$x_1^{(t)}$代入$x_2^{(t)}=\rho x_1^{(t)}+\sigma\varepsilon_{2t}$。}
\SolPlan 保留每轮两个中间值，完成两轮后把结果代回两条残差恒等式，专门检查第二坐标没有误用旧$x_1$。
\SolDerive
\textbf{第一轮。}得到
\[
\begin{aligned}
x_1^{(1)}&=\rho x_2^{(0)}+\sigma\varepsilon_{11}
=-\frac12,\\
x_2^{(1)}&=\rho x_1^{(1)}+\sigma\varepsilon_{21}
=-\frac14+\frac{\sqrt3}{2}.
\end{aligned}
\]
\textbf{第二轮。}得到
\[
\begin{aligned}
x_1^{(2)}&=\rho x_2^{(1)}+\sigma\varepsilon_{12}
=-\frac18-\frac{\sqrt3}{4},\\
x_2^{(2)}&=\rho x_1^{(2)}+\sigma\varepsilon_{22}
=-\frac1{16}-\frac{\sqrt3}{8}.
\end{aligned}
\]
\SolCheck 四个更新分别满足$x_1^{(t)}-\rho x_2^{(t-1)}=\sigma\varepsilon_{1t}$与$x_2^{(t)}-\rho x_1^{(t)}=\sigma\varepsilon_{2t}$；第二式确实使用同轮新$x_1^{(t)}$，因此可检出误用旧值的并行更新。
\SolAnswer 两轮轮末状态分别为$(-1/2,-1/4+\sqrt3/2)$与$(-1/8-\sqrt3/4,-1/16-\sqrt3/8)$。该确定性测试只验证更新次序与代数，不证明随机链已经收敛。
\end{solution}

% KNOWLEDGE-PLAN: KN-V5-C33-CONCEPT-001 L1
\paragraph{读前自检：抽样计算。}Gibbs一轮依次调用$d$个满条件抽样器；若第$j$个求值与抽样成本为$C_j,U_j$，请求和得到$T$轮时间$O(T\sum_j(C_j+U_j))$，并核对状态空间为$O(d)$。\par

\SLAdvancedSection{收敛、抽样计算与练习}{sec:V5-C04-S04}
\knowledgeanchor{SLM-19-05-006}{抽样计算}

概率图模型常把目标写成局部因子的乘积。以观测$y$、超参数$\alpha$、模型参数$\theta$和隐变量$z$为例，
\begin{equation}\label{eq:V5-C04-bayes-factorization}
\pi(\alpha,\theta,z\mid y)
\propto p(z,y\mid\theta)p(\theta\mid\alpha)p(\alpha).
\end{equation}
因此更新某个变量时，只需保留含该变量的因子。例如
\begin{align}
\pi(\alpha_i\mid\alpha_{-i},\theta,z,y)
&\propto p(\theta\mid\alpha)p(\alpha),\\
\pi(\theta_j\mid\theta_{-j},\alpha,z,y)
&\propto p(z,y\mid\theta)p(\theta\mid\alpha),\\
\pi(z_k\mid z_{-k},\alpha,\theta,y)
&\propto p(z,y\mid\theta).
\end{align}
图\ref{fig:V5-C04-bayes-markov-blanket}把这条“只看Markov毯”的消因子规则画成概率图；未写入当前条件核的每个因子都必须确实与当前变量无关。
% v2.6.0 figure UID: FIG-P641-01
% Image-guided reverse redraw: preserve exact factor-graph topology and isolate all annotations from paths.
\tikzset{slfig-FIG-P641-01/.style={font=\fontsize{9.2pt}{11.0pt}\selectfont,every path/.append style={line cap=round,line join=round}}}
\pgfkeysifdefined{/tikz/alt}{}{\tikzset{alt/.code={}}}
\begin{figure}[htbp]
\centering
\begin{tikzpicture}[slfig-FIG-P641-01,
  >=Stealth,
  var/.style={circle,draw=SLRuleGray,fill=white,minimum size=12mm,font=\fontsize{9.5pt}{11.4pt}\selectfont},
  blanket/.style={circle,draw=SLMidBlue,fill=SLSoftBlue,minimum size=12mm,font=\fontsize{9.5pt}{11.4pt}\selectfont},
  focus/.style={circle,draw=SLDeepBlue,fill=SLDeepBlue,text=white,
    minimum size=13mm,font=\fontsize{9.5pt}{11.4pt}\selectfont\bfseries},
  fac/.style={rectangle,draw=SLRuleGray,fill=SLSoftGray,minimum size=6mm,
    inner xsep=4pt,inner ysep=2.5pt,font=\fontsize{9.5pt}{11.4pt}\selectfont},
  activefac/.style={rectangle,draw=SLMidBlue,fill=SLSoftBlue,minimum size=7mm,
    inner xsep=4pt,inner ysep=2.5pt,font=\fontsize{9.5pt}{11.4pt}\selectfont},
  edge/.style={draw=SLGray,line width=.65pt},
  activeedge/.style={draw=SLBlue,line width=.95pt},
  scale=1.100,
  >={Stealth[length=2.3mm,width=1.6mm]},
  every path/.append style={line cap=round,line join=round},
  every node/.append style={line width=0.8pt},
  alt={对象—关系—结论：因子图中更新θ只需读取与θ相连的两个因子p(θ给定α)和p(z,y给定θ)；Markov毯变量为α,z,y，而与θ无关的因子p(α)可从满条件核中消去}]
\node[fac] (fa) at (-5.0,0) {$p(\alpha)$};
\node[blanket] (alpha) at (-3.5,0) {$\alpha$};
\node[activefac] (ft) at (-1.8,0) {$p(\theta\mid\alpha)$};
\node[focus] (theta) at (0,0) {$\theta$};
\node[activefac] (fzy) at (2.0,0) {$p(z,y\mid\theta)$};
\node[blanket] (z) at (4.0,1.05) {$z$};
\node[blanket] (y) at (4.0,-1.05) {$y$};

\draw[edge] (fa)--(alpha);
\draw[activeedge] (alpha)--(ft)--(theta)--(fzy)--(z);
\draw[activeedge] (fzy)--(y);

\foreach \v in {alpha,z,y}{
  \node[draw=SLTeal,line width=.70pt,dash pattern=on 2.2pt off 1.6pt,
    rounded corners=7pt,fit=(\v),inner sep=3pt] {};
}
\node[font=\fontsize{9.2pt}{11.0pt}\selectfont,text=SLTeal] at (1.0,1.80)
 {虚线圈：Markov 毯变量 $\alpha,z,y$};
\node[align=left,font=\fontsize{9.5pt}{11.4pt}\selectfont] at (0,-2.15)
 {$\displaystyle \pi(\theta\mid\alpha,z,y)\propto
 p(\theta\mid\alpha)\,p(z,y\mid\theta)$};
\node[align=center,font=\fontsize{9.5pt}{11.4pt}\selectfont,text=SLTextGray] at (-4.65,-1.48)
 {$p(\alpha)$与$\theta$无关\\置于毯变量标记外并消去};
\draw[line width=0.8pt,->,draw=SLTextGray] (-4.75,-.95) -- (fa.south);
\end{tikzpicture}
\captionsetup{width=.94\linewidth}
\caption{因子图中更新$\theta$只需读取与$\theta$相连的两个因子$p(\theta\mid\alpha)$和$p(z,y\mid\theta)$；Markov毯变量为$\alpha,z,y$，而与$\theta$无关的因子$p(\alpha)$可从满条件核中消去}
\label{fig:V5-C04-bayes-markov-blanket}
\end{figure}

% v2.6 layout: former forced clear removed; natural pagination.

\phantomsection\label{struct:V5-C04-CH19}
\textbf{伪代码。}算法\ref{alg:V5-C04-systematic-gibbs}把一次完整系统扫描定义为依次完成$d$个坐标更新。它保存轮末轨迹；若还保存轮内状态，必须使用独立下标，不把一次坐标更新误记成一整轮。

\phantomsection\label{struct:V5-C04-CH20}
\textbf{输入与输出。}输入为可归一化的目标、维数$d$、合法初值$x^{(0)}$、$d$个精确满条件抽样方法、彼此独立的随机数、燃烧轮数$m$、保留轮数$n$和可选函数$f$。输出包括完整轮末轨迹、逐坐标抽样记录、保留样本和可选样本均值；若某坐标更新失败，另行说明失败位置与本轮已经更新的中间状态。

\phantomsection\label{struct:V5-C04-CH21}
\textbf{参数含义。}$T=m+n$是完整扫描轮数；一轮含$d$次坐标抽样。$m$只决定哪些轮末状态不进入估计，不能证明链已收敛；$n=0$时允许只测试更新步骤，但均值、MCSE和ESS未定义。分块Gibbs把一个“坐标”替换为变量块，所有定义相应使用块满条件分布。

\phantomsection\label{struct:V5-C04-CH22}
\textbf{初始化。}开始前验证$\pi$可归一化、$x^{(0)}$位于目标支持、每个条件抽样器与式\eqref{eq:V5-C04-full-conditional}一致、随机源类型正确以及$m,n,d$合法。不能从支持外启动后期待退化条件分布自动把状态修回支持。

\phantomsection\label{struct:V5-C04-CH23}
\textbf{参数更新。}第$t$轮第$j$步使用当前轮内状态的$x_{-j}$调用$\mathcal C_j$，验证返回值属于$\mathcal X_j$且新完整状态位于目标支持，再立即覆盖$x_j$。后续坐标必须看到此前已更新的值。若某步失败，记录位置$(t,j)$；本轮中间状态不计作轮末样本。

\phantomsection\label{struct:V5-C04-CH24}
\textbf{停止条件。}正常分支恰好完成$T$轮；$T=0$时不调用随机数，返回只含初值的轨迹和空保留集。若条件分布未定义、抽样结果不是有限实数、样本越界、新状态离开支持或$f$不是有限实数，则立即停止。主轨迹只保留此前完成的轮末状态，失败轮的中间状态另行说明。

\phantomsection\label{struct:V5-C04-CH25}
\textbf{正确性与收敛条件。}定理\ref{thm:V5-C04-coordinate-reversible}与命题\ref{prop:V5-C04-scan-invariance}只证明目标保持。有限全支持时，定理\ref{thm:V5-C04-finite-positive}进一步给出时间平均；非周期性另用于固定时刻边缘收敛。一般受约束或连续状态必须按定义\ref{def:V5-C03-harris-conditions}和命题\ref{prop:V5-C03-general-convergence}分别核验正常Harris返、非周期性以及CLT/批均值所需的额外矩与混合条件。图\ref{fig:V5-C04-mixing-rho-comparison}说明即使这些结论成立，$|\rho|$接近$1$时统计效率仍可很低。

\phantomsection\label{struct:V5-C04-CH26}
\textbf{时间复杂度。}设第$j$个满条件抽样与支持验证的实际成本为$C_j$，一次坐标就地更新成本为$U_j$。完成$T=m+n$轮的时间为
\[
O\!\left(T\sum_{j=1}^d(C_j+U_j)\right).
\]
若每次轮末复制$d$维状态，还需$O(Td)$复制成本。概率图实现中$C_j$应按Markov毯因子求值计入，不能无条件写成$O(1)$。

\phantomsection\label{struct:V5-C04-CH27}
\textbf{空间复杂度。}只保留当前状态需要$O(d)$空间；保存$n$个轮末样本需要$O(nd)$；保存全部$T$轮及每次坐标更新可达$O(Td^2)$个标量。流式计算均值可把统计量附加空间降为$O(1)$或$O(\dim f)$，但用于诊断的轨迹不能在未记录替代信息时丢弃。

\phantomsection\label{struct:V5-C04-CH28}
\textbf{适用条件。}Gibbs抽样适合联合目标难以直接抽样、各满条件分布可精确且高效抽样、支持可由所选坐标或分块更新连通的模型。共轭层次模型、离散潜变量模型和稀疏概率图常满足这一结构。若只能近似抽满条件，应明确改用Metropolis-within-Gibbs、slice-within-Gibbs或其他带正确性证明的内核。

\pagebreak[4]
\phantomsection\label{struct:V5-C04-CH29}
\textbf{局限性。}强相关会使轴向路径混合缓慢；确定性约束会令单坐标条件分布退化；不当参数化可造成近不可约行为；高维全扫描成本可能很大；系统扫描缺少一般可逆性。分块、重参数化和随机扫描可能改善效率，但每次修改都必须重新核验目标保持与可达性。

\phantomsection\label{struct:V5-C04-CH30}
\textbf{常见错误。}不得把本轮尚未更新的坐标误写成新值；不得在条件归一化常数发散时继续抽样；不得把条件众数当条件随机样本；不得因“接受率为1”便跳过支持检查；不得把一次坐标更新与一次完整扫描混用时间下标；不得删除重复状态、只看单链轨迹或用固定燃烧期代替收敛条件。

\phantomsection\label{struct:V5-C04-CH31}
\textbf{易混淆概念对比。}
\begin{center}
\small
\renewcommand{\arraystretch}{1.60}
\begin{tabularx}{\linewidth}{@{}l X X@{}}
\toprule 对比 & Gibbs抽样 & 相关但不同的方法\\ \midrule
精确满条件与MH & 条件提议在正流分支接受率为1 & 其他提议需计算接受率并可能拒绝\\
系统扫描与随机扫描 & 固定顺序乘积核保持$\pi$，一般不再可逆 & 固定权重混合核可逆，但每步只更新一个坐标\\
条件抽样与条件众数 & 产生保持完整分布的随机链 & 条件众数更新是点优化，不给后验样本\\
单坐标与分块 & 单步便宜但可能强相关 & 分块成本较高，可能沿相关方向移动更远\\
正确性与效率 & 平稳性、可达性和返性决定正确性 & ESS、MCSE与运行时间描述给定正确内核的效率\\
\bottomrule
\end{tabularx}
\end{center}

\textbf{循环变量与更新顺序。}外层循环变量为完整扫描轮$t=1,\ldots,T$，内层循环变量为坐标$j=1,\ldots,d$；每次抽样成功后立即覆盖$x_j$，因此第$j+1$步读取本轮最新状态。

\textbf{判断条件。}每个坐标步核验条件分布是否定义、抽样值是否有限且属于坐标空间、新完整状态是否落在目标支持；可选函数$f$另行核验有限性。

\textbf{返回结果、异常与退化。}正常分支给出完整轮末轨迹和保留样本；中途停止时给出此前完整轮、$(t,j)$位置及本轮中间状态。$T=0$不调用随机数，$n=0$不构造样本均值。

\textbf{正确性依据。}定理\ref{thm:V5-C04-coordinate-reversible}保证每个坐标核保持$\pi$，命题\ref{prop:V5-C04-scan-invariance}保证其系统乘积仍保持$\pi$。

\textbf{不收敛或不适用情形。}目标不可归一化、满条件不相容或不可精确抽样、单坐标支持可约，或一般状态链不满足定义\ref{def:V5-C03-harris-conditions}中的正常Harris返条件时，本算法不能给出所需遍历均值保证；增加燃烧轮数不能修复这些问题。即便时间平均成立，固定时刻边缘收敛仍需非周期性，MCSE与批均值仍需命题\ref{prop:V5-C03-general-convergence}第3项的额外条件。

\textbf{复杂度假设。}复杂度使用本节定义的$C_j,U_j$，并显式计入状态复制；不把满条件归一化、因子求值、随机采样或支持验证无条件视为常数时间。

% ALGORITHM-FIELD-NOT-APPLICABLE: alg:V5-C04-systematic-gibbs | 训练时间复杂度=本算法从给定目标分布抽样，不拟合训练参数; 预测时间复杂度=本算法不单独执行预测，成本按每轮满条件抽样报告
\begingroup
\SetArgSty{textnormal}
% KNOWLEDGE-PLAN: KN-V5-C33-ALGORITHM_IDEA-003 L2

% KNOWLEDGE-PLAN: KN-V5-C33-ALGORITHM_IDEA-004 L2
\begin{keypointbox}[title={带部分扫描诊断的系统Gibbs抽样：五步阅读}]
\textbf{输入。}写出数据对象、固定参数与必须成立的条件。\par
\textbf{初始化。}标明主状态、计数器和第一个合法状态。\par
\textbf{核心更新。}只手算一次从旧状态到候选状态的更新，并指出何时提交。\par
\textbf{数学停止。}写出能直接核验的停止证书；预算耗尽不等同于收敛。\par
\textbf{输出。}列出返回对象、实际计数、中心状态和用于复核的诊断。
\end{keypointbox}

\begin{AlgorithmContract}{
  input={可归一化目标、$d$个精确满条件抽样接口、合法初值、独立随机源、燃烧/保留扫描数$m,n$及可选$f$},
  output={最后合法轮末轨迹、只读部分扫描状态、保留样本、完整扫描$t_{\rm done}$、成功坐标抽样$c_{\rm done}$、随机调用$q_{\rm done}$、状态与最终诊断},
  preconditions={$d\ge1,m,n\in\mathbb N_0$；目标与满条件相容；初值在支持内；各条件抽样接口/坐标空间与随机源契约合法},
  initialization={$T=m+n,x=x^{(0)}$，$t_{\rm done}=c_{\rm done}=q_{\rm done}=0$，轮末轨迹只含初值},
  loop={按固定顺序逐扫描、逐坐标从当前最新工作状态的满条件抽样},
  update={单坐标样本及新完整状态合法后才原子提交到工作状态；整轮完成后才追加轮末轨迹并增加$t_{\rm done}$},
  domain={每个工作状态在目标支持；条件样本有限且属于对应坐标空间；轮末轨迹只含完整扫描},
  failure={接口/初值/预算非法返回$\mathtt{invalid\_input}$；条件随机源失败返回$\mathtt{random\_source\_failure}$；条件归一化、新状态或$f$求值失败返回$\mathtt{numerical\_failure}$，主轨迹回滚到上一完整扫描},
  stop={完成$T$次扫描、$T=0$或首次坐标失败时停止；固定扫描完成不称收敛},
  budget={至多$T d$次满条件随机调用及相应支持检查},
  status={$\mathtt{completed}$、$\mathtt{invalid\_input}$、$\mathtt{numerical\_failure}$、$\mathtt{random\_source\_failure}$},
  iterations={$t_{\rm done}$计完整提交扫描，$c_{\rm done}$计成功坐标更新，$q_{\rm done}$计条件随机源调用},
  diagnostics={失败$(t,j,\mathrm{stage})$、只读部分扫描状态、各坐标工作量、保留段/均值状态和计划/实际扫描},
  complexity={$\Theta(T\sum_j(C_j+U_j+d))$时间；轮末轨迹$\Theta(Td)$空间，工作状态$\Theta(d)$}
}
% v2.6 layout: former forced clear removed; natural pagination.
\begin{algorithm}[tbp]
\caption{带部分扫描诊断的系统Gibbs抽样}\label{alg:V5-C04-systematic-gibbs}
\KwIn{可归一化的目标；$d$个精确满条件抽样方法$\mathcal C_j$；合法$x^{(0)}$；独立随机数；$m,n$；可选$f$}
\KwOut{最后合法轮末轨迹、部分扫描/保留段；$t_{\rm done},c_{\rm done},q_{\rm done}$；$\mathtt{status}$与最终诊断}
若$d,m,n$、目标/初值、满条件接口或随机源契约非法，则返回空结果、计数$0$、$\mathtt{invalid\_input}$及字段诊断\;
令$T=m+n,x=x^{(0)},t_{\rm done}=c_{\rm done}=q_{\rm done}=0$，轮末轨迹$\mathcal T=(x^{(0)})$；若$T=0$，返回该轨迹、空保留段、$\mathtt{completed}$及零预算诊断\;
\For{$t\leftarrow1,\ldots,T$}{
  令$x^{\mathrm{start}}\leftarrow x,x_{\rm work}\leftarrow x$，本轮坐标记录为空\;
  \For{$j\leftarrow1,\ldots,d$}{
    使用$x_{{\rm work},-j}$调用$\mathcal C_j$并增加$q_{\rm done}$；调用失败则返回前$t-1$轮、只读$x_{\rm work}$、$\mathtt{random\_source\_failure}$及$(t,j)$诊断\;
    若条件归一化诊断、$y_j$有限性/坐标域或新完整状态支持检查失败，则返回上一完整轨迹、只读部分状态、$\mathtt{numerical\_failure}$及$(t,j)$诊断\;
    原子令$x_{{\rm work},j}\leftarrow y_j$，增加$c_{\rm done}$并记录条件上下文/轮内状态\;
  }
  原子提交$x\leftarrow x_{\rm work}$并向$\mathcal T$追加；令$t_{\rm done}\leftarrow t$，记录本轮变化\;
}
令保留集$\mathcal S\leftarrow(X^{(m+1)},\ldots,X^{(m+n)})$\;
若给出$f$且$n>0$，核验所有$f(X)$有限并计算均值；失败时保留轨迹/样本并返回$\mathtt{numerical\_failure}$及均值阶段诊断\;
返回$\mathcal T,\mathcal S$、实际计数、$\mathtt{completed}$及扫描/保留段最终诊断；$n=0$时均值未定义\;
\end{algorithm}
\end{AlgorithmContract}
\endgroup

\phantomsection\label{struct:V5-C04-CH32}
\Needspace{6\baselineskip}
\begin{example}[五分类模型的可复算Gibbs目标]\label{exm:V5-C04-five-category}
五类概率为
\[
p_1=\frac\theta4+\frac18,
\quad p_2=\frac\theta4,
\quad p_3=\frac\eta4,
\quad p_4=\frac\eta4+\frac38,
\quad p_5=\frac{1-\theta-\eta}{2},
\]
参数域为$D=\{(\theta,\eta):\theta>0,\eta>0,\theta+\eta<1\}$，观测计数为$(14,1,1,1,5)$。仅有似然不足以定义后验；本例明确补充$D$上的均匀先验。任务是推导后验核、两个满条件分布的精确有限混合表示，并给出无需随机运行的精确矩基准。
\end{example}

\phantomsection\label{struct:V5-C04-CH33}
\Needspace{4\baselineskip}
\SLExampleSolutionHeading{exm:V5-C04-five-category}
\begin{solution}
\SLReadTranslation{用$D$上的均匀先验与多项计数似然定义三角域后验，再把两个满条件写成可直接抽样的有限Beta混合，并给出不依赖随机运行的精确矩基准。}
\SolGiven 参数域$D$有界，先验在$D$上均匀；观测计数固定，且$\theta,\eta,1-\theta-\eta$在$D$内均为正。
\SLMethodTrigger{先写后验核；固定一个参数后把另一个缩放到$(0,1)$并作有限多项式展开，从而识别Beta混合；最后用单纯形积分计算矩。}
\SolPlan 依次得到后验核、$u=\theta/(1-\eta)$的15分量Beta混合、$v=\eta/(1-\theta)$的二分量Beta混合，再用Dirichlet积分恒等式求四个后验矩。
\SolDerive
\textbf{后验核。}多项式系数以及$8,4,2$的幂与参数无关，可以并入归一化常数。因此
\begin{equation}\label{eq:V5-C04-five-posterior}
\pi(\theta,\eta\mid y)
\propto
(1+2\theta)^{14}\theta\eta(3+2\eta)
(1-\theta-\eta)^5\mathbf 1_D(\theta,\eta).
\end{equation}
\textbf{$\theta$满条件。}固定$\eta$并令$s=1-\eta$、$u=\theta/s$。二项展开给出
\[
(1+2su)^{14}u(1-u)^5
=\sum_{r=0}^{14}\binom{14}{r}(2s)^r
u^{r+1}(1-u)^5.
\]
所以$u\mid\eta$是$\operatorname{Beta}(r+2,6)$的15分量混合，分量权重为
\begin{equation}\label{eq:V5-C04-theta-mixture}
\omega_r^{(\theta)}(s)
=\frac{\binom{14}{r}(2s)^r B(r+2,6)}
{\sum_{\ell=0}^{14}\binom{14}{\ell}(2s)^\ell B(\ell+2,6)},
\quad r=0,\ldots,14,
\end{equation}
抽得$u$后令$\theta=su$。

\textbf{$\eta$满条件。}固定$\theta$并令$s=1-\theta$、$v=\eta/s$，则
\[
\eta(3+2\eta)(s-\eta)^5\,\mathrm d\eta
\propto
\bigl[3v(1-v)^5+2s v^2(1-v)^5\bigr],\mathrm dv.
\]
因此$v\mid\theta$是$\operatorname{Beta}(2,6)$和$\operatorname{Beta}(3,6)$的二分量混合，未归一化权重分别为$3B(2,6)$与$2sB(3,6)$；抽得$v$后令$\eta=sv$。两步都保持$D$，且是精确满条件抽样。

\textbf{确定性矩基准。}定义
\[
I_{a,b}=\int_D \theta^a\eta^b
(1+2\theta)^{14}\theta\eta(3+2\eta)
(1-\theta-\eta)^5\,\mathrm d\theta\,\mathrm d\eta.
\]
计算时先用二项式展开$(1+2\theta)^{14}$，再把$3+2\eta$拆成两项。对$A,B,C>0$，所用的单纯形Dirichlet恒等式为
\[
\int_D\theta^{A-1}\eta^{B-1}(1-\theta-\eta)^{C-1}
\,\mathrm d\theta\,\mathrm d\eta
=\frac{\Gamma(A)\Gamma(B)\Gamma(C)}{\Gamma(A+B+C)}.
\]
对展开式的第$r$项，常数$3$对应
$(A,B,C)=(r+a+2,b+2,6)$，而$2\eta$对应
$(A,B,C)=(r+a+2,b+3,6)$。逐项代入即得
\begin{align}
I_{a,b}
=\sum_{r=0}^{14}\binom{14}{r}2^r\Bigg[
&\frac{3\,\Gamma(r+a+2)\Gamma(b+2)\Gamma(6)}
{\Gamma(r+a+b+10)}\\
&+\frac{2\,\Gamma(r+a+2)\Gamma(b+3)\Gamma(6)}
{\Gamma(r+a+b+11)}\Bigg].
\end{align}
于是$E(\theta^a\eta^b\mid y)=I_{a,b}/I_{0,0}$。确定性复算结果为
\begin{align*}
E(\theta\mid y)&\approx0.51995509,
&\operatorname{Var}(\theta\mid y)&\approx0.01776292,\\
E(\eta\mid y)&\approx0.12316991,
&\operatorname{Var}(\eta\mid y)&\approx0.00655205.
\end{align*}
\SolCheck 后验核在有界域$D$上非负可积；两组混合权重归一化后均和为$1$。$u,v\in(0,1)$分别保证$0<\theta<1-\eta$与$0<\eta<1-\theta$。独立的三角域数值积分复得上述四个矩，并给出$E(\theta+\eta\mid y)\approx0.64312500<1$。

\SolAnswer 令$u=\theta/(1-\eta)$，则$u\mid\eta,y$是15分量$\operatorname{Beta}(r+2,6)$混合，且$\theta=(1-\eta)u$；令$v=\eta/(1-\theta)$，则$v\mid\theta,y$是$\operatorname{Beta}(2,6)$与$\operatorname{Beta}(3,6)$混合，且$\eta=(1-\theta)v$。后验均值约为$0.51995509,0.12316991$，方差约为$0.01776292,0.00655205$。这些数值是确定性积分基准，不是一次随机Gibbs运行的样本结果。
\end{solution}

\phantomsection\label{struct:V5-C04-CH34}
\begin{chapterexercisebox}
\phantomsection\label{struct:V5-C04-CH35}
本章共12道练习。每道题干后立即给出同号解析；长解析可自然跨页，续页标题保留题号。
\end{chapterexercisebox}

\begin{SLChapterExercisePairSection}

\SLSourceBookExercises

\Needspace{6\baselineskip}
% EXERCISE-SOURCE: original_book; source_id=EXR-19-0008; source_number=19.8
% SOLUTION-SOURCE: original_book; source_id=EXR-19-0008; source_number=19.8
\begin{exercise}\label{ex:V5-C04-S04-01}
五分类概率、参数域和观测计数如例\ref{exm:V5-C04-five-category}。回答：（a）为什么只给似然不能唯一确定“参数均值与方差”；（b）在$D$上均匀先验下推导后验核；（c）推导$\theta\mid\eta,y$与$\eta\mid\theta,y$的精确有限Beta混合抽样法；（d）给出后验均值、方差的可复算确定性基准；（e）写明一次真实Gibbs估计还必须保存哪些证据。不得在没有随机流时虚构样本均值。
\end{exercise}
\ifSLFullEdition
\phantomsection\label{sol:V5-C04-S04-01}
\begin{chapterexercisesolution}{ex:V5-C04-S04-01}
（a）参数的“均值与方差”必须相对于一个概率分布定义。似然$L(\theta,\eta;y)$只描述给定参数时数据的相对可能性；不同先验会给出不同后验矩。因此原题数据还需先验或其他目标分布约定。

（b）在$D$上取均匀先验后，多项式似然中与参数无关的常数消去，得到
\[
\pi(\theta,\eta\mid y)
\propto(1+2\theta)^{14}\theta\eta(3+2\eta)
(1-\theta-\eta)^5\mathbf1_D,
\]
与式\eqref{eq:V5-C04-five-posterior}一致。

（c）固定$\eta$，令$s=1-\eta$、$u=\theta/s$。展开$(1+2su)^{14}$后，$u$是$\operatorname{Beta}(r+2,6)$、$r=0,\ldots,14$的混合，权重正是式\eqref{eq:V5-C04-theta-mixture}。固定$\theta$，令$s=1-\theta$、$v=\eta/s$，则$v$的核为
\[
3v(1-v)^5+2s v^2(1-v)^5,
\]
所以两分量分别为$\operatorname{Beta}(2,6)$与$\operatorname{Beta}(3,6)$，未归一化权重为$3B(2,6)$与$2sB(3,6)$。每次抽样后分别令$\theta=su$、$\eta=sv$，状态始终留在$D$。

（d）对例\ref{exm:V5-C04-five-category}中的$I_{a,b}$求比值得
\[
E(\theta\mid y)\approx 0.51995509,
\quad \operatorname{Var}(\theta\mid y)\approx 0.01776292,
\]
\[
E(\eta\mid y)\approx 0.12316991,
\quad \operatorname{Var}(\eta\mid y)\approx 0.00655205.
\]
这些数来自有限Gamma函数和，不依赖随机运行。

（e）真实报告至少说明先验、初值、扫描顺序、混合分量、随机数生成原则、完整轮末轨迹、燃烧与保留规则、MCSE/ESS，以及样本矩对上述精确矩的误差。没有实际运行时只能给出算法步骤与解析基准，不能报告一次虚构的Gibbs估计。
\end{chapterexercisesolution}
\fi

\SLLectureSupplementExercises

\Needspace{6\baselineskip}
% EXERCISE-SOURCE: lecture_supplement
% SOLUTION-SOURCE: lecture_supplement
\begin{exercise}\label{ex:V5-C04-S01-01}
设$(X_1,X_2)\in\{0,1\}^2$的联合质量依次为
\[
\pi(0,0)=\frac1{10},\quad
\pi(0,1)=\frac2{10},\quad
\pi(1,0)=\frac3{10},\quad
\pi(1,1)=\frac4{10}.
\]
计算两个坐标的全部满条件概率，并从$(0,0)$出发计算一轮“先$X_1$后$X_2$”系统扫描的四个轮末转移概率。
\end{exercise}
\ifSLFullEdition
\phantomsection\label{sol:V5-C04-S01-01}
\begin{chapterexercisesolution}{ex:V5-C04-S01-01}
固定$X_2=0$时，该列总质量为$(1+3)/10$，所以
\[
P(X_1=1\mid X_2=0)=\frac34,
\quad P(X_1=0\mid X_2=0)=\frac14.
\]
固定$X_2=1$得到$P(X_1=1\mid X_2=1)=4/(2+4)=2/3$。按行归一化又有
\[
P(X_2=1\mid X_1=0)=\frac23,
\qquad
P(X_2=1\mid X_1=1)=\frac47.
\]
从$(0,0)$先更新$X_1$。若新$X_1=0$，概率为$1/4$，随后$X_2$取$0,1$的概率为$1/3,2/3$；若新$X_1=1$，概率为$3/4$，随后概率为$3/7,4/7$。因此
\[
K_{12}((0,0),\cdot)
=\left(\frac1{12},\frac16,\frac9{28},\frac37\right)
\]
按$(0,0),(0,1),(1,0),(1,1)$排列。四项之和为$1$。
\end{chapterexercisesolution}
\fi

\Needspace{6\baselineskip}
% EXERCISE-SOURCE: lecture_supplement
% SOLUTION-SOURCE: lecture_supplement
\begin{exercise}\label{ex:V5-C04-S01-02}
目标密度在三角域$D=\{(x_1,x_2):x_1>0,x_2>0,x_1+x_2<1\}$上为常数$2$，域外为$0$。推导两个满条件密度，说明它们的支持怎样随另一坐标改变，并解释为什么不能在$x_2=1$处套用密度比定义$x_1\mid x_2$。
\end{exercise}
\ifSLFullEdition
\phantomsection\label{sol:V5-C04-S01-02}
\begin{chapterexercisesolution}{ex:V5-C04-S01-02}
对$0<x_2<1$，边缘密度为
\[
\pi_2(x_2)=\int_0^{1-x_2}2\,\mathrm dx_1=2(1-x_2).
\]
故两个满条件概率核为
\[
\pi_1(x_1\mid x_2)
=\frac{\boldsymbol1_{(0,\,1-x_2)}(x_1)}{1-x_2},
\qquad
\pi_2(x_2\mid x_1)
=\frac{\boldsymbol1_{(0,\,1-x_1)}(x_2)}{1-x_1}.
\]
各自积分为$1$，且另一坐标越大，条件支持越短。$x_2=1$时
\[
\pi_2(1)=\int_0^0 2\,\mathrm dx_1=0,\qquad
\frac{\pi(x_1,1)}{\pi_2(1)}=\frac00,
\]
该条件事件在目标下为零概率；只能选取正则条件分布的一个版本，不能由密度比唯一指定。
\end{chapterexercisesolution}
\fi

\Needspace{6\baselineskip}
% EXERCISE-SOURCE: lecture_supplement
% SOLUTION-SOURCE: lecture_supplement
\begin{exercise}\label{ex:V5-C04-S02-01}
目标在$\{(0,0),(1,1)\}$上各赋概率$1/2$。写出两个满条件分布和系统Gibbs核，判断平稳性、可逆性、不可约性与唯一平稳分布。给出一种能连接两个状态的合法分块更新。
\end{exercise}
\ifSLFullEdition
\phantomsection\label{sol:V5-C04-S02-01}
\begin{chapterexercisesolution}{ex:V5-C04-S02-01}
若$X_2=0$，目标支持只允许$X_1=0$；若$X_2=1$，只允许$X_1=1$。因此
\[
P(X_1=x_2\mid X_2=x_2)=1,\qquad
P(X_2=x_1\mid X_1=x_1)=1.
\]
任何单坐标更新都保持原状态，在状态顺序$(0,0),(1,1)$下，
\[
K_{\rm sys}=I_2,\qquad
\pi K_{\rm sys}=\pi,\qquad
\pi_iK_{ij}=\pi_jK_{ji}.
\]
所以它关于均匀目标可逆并保持目标，但可约；事实上任意两个状态上的概率分布都平稳，故不唯一。

把两个坐标作为一个块并从联合目标抽样，得到
\[
K_{\rm block}
=\begin{pmatrix}1/2&1/2\\1/2&1/2\end{pmatrix},
\qquad \pi K_{\rm block}=\pi.
\]
该核一步连接两个状态。
\end{chapterexercisesolution}
\fi

\Needspace{6\baselineskip}
% EXERCISE-SOURCE: lecture_supplement
% SOLUTION-SOURCE: lecture_supplement
\begin{exercise}\label{ex:V5-C04-S02-02}
在$\{-1,1\}^3$上令
\[
\pi(x)\propto\exp\{\beta(x_1x_2+x_2x_3)\}.
\]
推导三个Bernoulli满条件概率。取$\beta=(\log3)/2$、初值$(-1,1,-1)$，按$1,2,3$更新，并用固定均匀数$(0.2,0.9,0.4)$执行一轮；约定$u\le P(X_j=1\mid x_{-j})$时取$1$。
\end{exercise}
\ifSLFullEdition
\phantomsection\label{sol:V5-C04-S02-02}
\begin{chapterexercisesolution}{ex:V5-C04-S02-02}
只保留含当前坐标的指数项。对$s\in\{-1,1\}$，
\[
P(X_1=s\mid x_2)
=\frac{e^{\beta s x_2}}{e^{\beta x_2}+e^{-\beta x_2}},
\quad
P(X_3=s\mid x_2)
=\frac{e^{\beta s x_2}}{e^{\beta x_2}+e^{-\beta x_2}},
\]
\[
P(X_2=s\mid x_1,x_3)
=\frac{e^{\beta s(x_1+x_3)}}
{e^{\beta(x_1+x_3)}+e^{-\beta(x_1+x_3)}}.
\]
当$\beta=(\log3)/2$且当前$x_2=1$时，$P(X_1=1\mid x_2)=3/4$；$u_1=0.2$给出$x_1=1$。此后$x_1+x_3=0$，故$P(X_2=1\mid x_1,x_3)=1/2$；$u_2=0.9$给出$x_2=-1$。最后$P(X_3=1\mid x_2=-1)=1/4$；$u_3=0.4$给出$x_3=-1$。轮末状态为$(1,-1,-1)$。这是固定输入下的精确执行结果。
\end{chapterexercisesolution}
\fi

\Needspace{6\baselineskip}
% EXERCISE-SOURCE: lecture_supplement
% SOLUTION-SOURCE: lecture_supplement
\begin{exercise}\label{ex:V5-C04-S03-01}
把精确满条件$\pi_j(\cdot\mid x_{-j})$代入单坐标\MetropolisHastings{}接受比，逐项证明正流分支的接受率为$1$。若改用一般提议$q_j(y_j\mid x_j,x_{-j})$，写出正确接受率，并说明反向提议为零时怎样处理。
\end{exercise}
\ifSLFullEdition
\phantomsection\label{sol:V5-C04-S03-01}
\begin{chapterexercisesolution}{ex:V5-C04-S03-01}
单坐标候选满足$y_{-j}=x_{-j}$。在$\pi(x)>0$、$\pi(y)>0$且正反条件密度为正的分支，
\[
\frac{\pi(y)\pi_j(x_j\mid x_{-j})}
{\pi(x)\pi_j(y_j\mid x_{-j})}=1
\]
的约分过程与正文相同，所以$\alpha=1$。候选由精确满条件产生时，目标零质量点的抽样概率为零；零边缘条件值则不能用该比值定义。

一般提议$q_j$下，正确接受率为
\[
\alpha(x,y)=\min\!\left\{1,
\frac{\widetilde\pi(y)
q_j(x_j\mid y_j,x_{-j})}
{\widetilde\pi(x)
q_j(y_j\mid x_j,x_{-j})}\right\}.
\]
若正向流为正而反向提议为零，则双向公共流为零，必须确定拒绝；不能用一个人为小正数替代零。
\end{chapterexercisesolution}
\fi

\Needspace{6\baselineskip}
% EXERCISE-SOURCE: lecture_supplement
% SOLUTION-SOURCE: lecture_supplement
\begin{exercise}\label{ex:V5-C04-S03-02}
对命题\ref{prop:V5-C04-bivariate-normal}中的目标取$\rho=0.8$。推导两个满条件分布、$X_2$轮末链的AR系数与创新方差，并计算理论积分自相关时间及渐近ESS比例。解释“每步接受”为什么不等于“样本近似独立”。
\end{exercise}
\ifSLFullEdition
\phantomsection\label{sol:V5-C04-S03-02}
\begin{chapterexercisesolution}{ex:V5-C04-S03-02}
$\rho=0.8=4/5$时，$1-\rho^2=9/25$，所以
\[
X_1\mid X_2=x_2\sim N\!\left(\frac45x_2,\frac9{25}\right),
\qquad
X_2\mid X_1=x_1\sim N\!\left(\frac45x_1,\frac9{25}\right).
\]
轮末$X_2$的AR系数为$\rho^2=16/25$，创新方差为$1-\rho^4=1-256/625=369/625$。理论积分自相关时间和ESS比例为
\[
\tau_{\mathrm{int}}
=\frac{1+16/25}{1-16/25}=\frac{41}{9}\approx4.5556,
\qquad
\frac{N_{\mathrm{eff}}}{N}=\frac9{41}\approx0.2195.
\]
每个条件候选都被采用只说明没有MH拒绝；相邻状态仍通过旧坐标传递相关性，因此远非独立。
\end{chapterexercisesolution}
\fi

\Needspace{6\baselineskip}
% EXERCISE-SOURCE: lecture_supplement
% SOLUTION-SOURCE: lecture_supplement
\begin{exercise}\label{ex:V5-C04-S03-03}
在同一二元正态目标下取$\rho=1/2$、$x^{(0)}=(0,0)$，两轮创新依次为$(1,0)$和$(0,-1)$。精确计算两轮轮末状态，并说明这些固定输入能检验什么、不能证明什么。
\end{exercise}
\ifSLFullEdition
\phantomsection\label{sol:V5-C04-S03-03}
\begin{chapterexercisesolution}{ex:V5-C04-S03-03}
此时$\sigma=\sqrt3/2$。第一轮
\[
x_1^{(1)}=\frac12\cdot0+\frac{\sqrt3}{2}\cdot1=\frac{\sqrt3}{2},
\qquad
x_2^{(1)}=\frac12x_1^{(1)}+\frac{\sqrt3}{2}\cdot0=\frac{\sqrt3}{4}.
\]
第二轮
\[
x_1^{(2)}=\frac12x_2^{(1)}=\frac{\sqrt3}{8},
\qquad
x_2^{(2)}=\frac12x_1^{(2)}-\frac{\sqrt3}{2}
=-\frac{7\sqrt3}{16}.
\]
这些值可检验更新顺序、旧值/新值使用、均值和标准差系数；两轮固定创新不能证明随机数生成器正确、链收敛或样本矩接近目标矩。
\end{chapterexercisesolution}
\fi

\Needspace{6\baselineskip}
% EXERCISE-SOURCE: lecture_supplement
% SOLUTION-SOURCE: lecture_supplement
\begin{exercise}\label{ex:V5-C04-S03-04}
证明固定权重随机扫描Gibbs核可逆。再使用练习\ref{ex:V5-C04-S01-01}的联合分布，分别计算从$(0,0)$到$(1,1)$的一轮正序扫描$K_1K_2$与反序扫描$K_2K_1$概率，据此说明系统扫描一般不具有可逆性。
\end{exercise}
\ifSLFullEdition
\phantomsection\label{sol:V5-C04-S03-04}
\begin{chapterexercisesolution}{ex:V5-C04-S03-04}
每个$K_j$关于$\pi$的双向流对称，常数加权和
$K_{\mathrm{rs}}=\sum_jw_jK_j$仍满足
\[
\pi(\mathrm dx)K_{\mathrm{rs}}(x,\mathrm dy)
=\pi(\mathrm dy)K_{\mathrm{rs}}(y,\mathrm dx),
\]
所以固定权重随机扫描可逆。

对练习\ref{ex:V5-C04-S01-01}，正序从$(0,0)$到$(1,1)$的概率为
\[
K_1K_2((0,0),(1,1))=\frac34\cdot\frac47=\frac37.
\]
反序则为
\[
K_2K_1((0,0),(1,1))=\frac23\cdot\frac23=\frac49.
\]
两者不等，故$K_1,K_2$不交换。由于$(K_1K_2)^*=K_2K_1$，正序系统扫描不等于自身伴随核，因而在这个例子中不具可逆性；两种顺序仍都保持$\pi$。
\end{chapterexercisesolution}
\fi

\Needspace{6\baselineskip}
% EXERCISE-SOURCE: lecture_supplement
% SOLUTION-SOURCE: lecture_supplement
\begin{exercise}\label{ex:V5-C04-S04-02}
设
\[
p(\alpha,\theta,z,y)
=p(\alpha)p(\theta\mid\alpha)
\prod_{k=1}^N p(z_k\mid\theta)p(y_k\mid z_k,\theta).
\]
分别写出$\alpha$、$\theta$和单个$z_k$的满条件分布中必须保留的因子，并解释为什么其他因子可以消去。若$z_k$只在一个因子中出现，说明其Markov毯。
\end{exercise}
\ifSLFullEdition
\phantomsection\label{sol:V5-C04-S04-02}
\begin{chapterexercisesolution}{ex:V5-C04-S04-02}
条件化时，与当前变量无关的因子并入归一化常数。于是
\[
p(\alpha\mid\theta,z,y)
\propto p(\alpha)p(\theta\mid\alpha),
\]
\[
p(\theta\mid\alpha,z,y)
\propto p(\theta\mid\alpha)
\prod_{k=1}^N p(z_k\mid\theta)p(y_k\mid z_k,\theta),
\]
而对单个$z_k$，
\[
p(z_k\mid z_{-k},\alpha,\theta,y)
\propto p(z_k\mid\theta)p(y_k\mid z_k,\theta).
\]
其他观测与隐变量因子不含$z_k$，所以对$z_k$归一化时消去。$z_k$的Markov毯由它的父变量$\theta$、子变量$y_k$以及$y_k$的其他父变量组成；在给定因子结构中没有额外变量。
\end{chapterexercisesolution}
\fi

\Needspace{6\baselineskip}
% EXERCISE-SOURCE: lecture_supplement
% SOLUTION-SOURCE: lecture_supplement
\begin{exercise}\label{ex:V5-C04-S04-03}
算法\ref{alg:V5-C04-systematic-gibbs}计划执行$T$轮、每轮$d$个坐标，第$t$轮第$j$个坐标失败。回答：主轮末轨迹应含多少个状态，完成轮数是多少，失败轮怎样返回；若每坐标成本上界为$C$且轮末复制成本为$O(d)$，失败前的时间上界和保存全部轮内状态的空间上界分别是什么？
\end{exercise}
\ifSLFullEdition
\phantomsection\label{sol:V5-C04-S04-03}
\begin{chapterexercisesolution}{ex:V5-C04-S04-03}
停止前已完成$t-1$轮，所以完整轮数为$t-1$；主轮末轨迹还含初值，共$t$个状态。第$t$轮只完成了前$j-1$个坐标更新，当前中间状态必须单列，不能追加为一个轮末样本。
\[
N_{\rm round}=t-1,\qquad
|\mathcal T_{\rm main}|=t,\qquad
N_{\rm coordinate}=(t-1)d+(j-1).
\]

失败前最多完成$(t-1)d+(j-1)$次坐标更新，坐标成本上界为
\[
O\bigl(((t-1)d+j-1)C\bigr).
\]
再计入$t-1$次轮末复制得到$O(((t-1)d+j)C+(t-1)d)$的参数化上界；若把更新成本统一并入$C$，可简写为$O(tdC)$。保存至多$T$轮的全部$d$个轮内$d$维状态需要$O(Td^2)$标量空间；只保存轮末状态则是$O(Td)$。
\[
S_{\rm all\ inner}=O(Td^2),\qquad
S_{\rm endpoints}=O(Td).
\]
\end{chapterexercisesolution}
\fi

\Needspace{6\baselineskip}
% EXERCISE-SOURCE: lecture_supplement
% SOLUTION-SOURCE: lecture_supplement
\begin{exercise}\label{ex:V5-C04-S04-04}
判断下列说法，并为错误说法给出反例或缺失条件：（a）任意给出的一组满条件分布都对应某个联合分布；（b）Gibbs每步接受，所以一定混合很快；（c）足够长的燃烧期可修复可约链；（d）系统扫描总是可逆；（e）任意依赖当前状态的坐标选择概率都保持目标分布；（f）漂亮的单链轨迹足以证明程序和目标都正确。
\end{exercise}
\ifSLFullEdition
\phantomsection\label{sol:V5-C04-S04-04}
\begin{chapterexercisesolution}{ex:V5-C04-S04-04}
六个说法都缺条件或错误。
\[
\begin{aligned}
P(X=1\mid Y=0)&=P(X=1\mid Y=1)=\frac12,\\
P(Y=1\mid X=0)&=\frac14,\qquad
P(Y=1\mid X=1)=\frac34
\end{aligned}
\]
不能来自同一严格正联合分布：前一行推出$X$与$Y$独立，后一行却要求依赖，给出（a）的具体不相容反例。
\begin{enumerate}[label=(\alph{enumi}),leftmargin=2.4em]
  \item 满条件必须彼此相容并对应同一个可归一化联合分布；任意指定的一组条件分布可能不存在共同联合分布。
  \item 练习\ref{ex:V5-C04-S03-02}中每步都采用新值，但$\rho=0.8$时理论ESS比例约为$0.2195$，所以接受不保证快速混合。
  \item 练习\ref{ex:V5-C04-S02-01}的核为恒等核；丢弃任意多步后仍停在初值。
  \item 练习\ref{ex:V5-C04-S03-04}给出$K_1K_2\ne K_2K_1$的具体系统扫描，因此一般不具可逆性。
  \item 状态依赖选择权重会改变正反更新流；除非另有平衡证明或校正，不能沿用固定权重混合核的证明。
  \item 单链轨迹只能显示已访问路径。错误目标、共同遗漏的模态或程序与理论同时犯错都可能产生平滑轨迹；还需目标、转移核、独立测试、多个初值和定量诊断证据。
\end{enumerate}
\[
\begin{aligned}
K=I&\Longrightarrow K^m=I
  &&(\text{燃烧期不修复可约性}),\\
(K_1K_2)^*&=K_2K_1\ne K_1K_2
  &&(\text{系统扫描一般不可逆}).
\end{aligned}
\]
\end{chapterexercisesolution}
\fi

\end{SLChapterExercisePairSection}
% Reserve the complete closing unit so a lone fifth checklist item cannot spill.
\Needspace{24\baselineskip}
\begin{SLCompactClosure}
\phantomsection\label{struct:V5-C04-CH36}
\SLChapterFiveSummary{满条件首先是可测概率核；正边缘处由密度比给出，零边缘集合上必须选定版本。}
{单坐标核的条件概率乘积关于新旧坐标对称，故可逆并保持$\pi$；系统扫描保持$\pi$但通常不可逆。}
{一轮依固定次序从当前轮内状态抽取各坐标，后一步立即读取前一步新值。}
{适合满条件可精确抽样且坐标移动连通目标支持的模型；强相关与退化约束会导致慢混合或可约。}
{\NextChapter{35}{潜在狄利克雷分配}，把本章坐标Gibbs用于主题后验推断。}

\phantomsection\label{struct:V5-C04-CH37}
\begin{selfcheckbox}
\small
\begin{enumerate}[leftmargin=2.2em,itemsep=0.15em,topsep=0.25em,parsep=0pt]
  \item 能否从联合密度得到式\eqref{eq:V5-C04-full-conditional}并处理零边缘质量，再用对称条件流证明单坐标核可逆、说明系统扫描只自动继承平稳性？
  \item 能否检查单坐标支持连通性、构造可约Gibbs反例，并完整执行算法\ref{alg:V5-C04-systematic-gibbs}以区分轮末轨迹与部分扫描？
  \item 能否推导二元正态满条件、AR递推、积分自相关时间与ESS比例，同时从概率图因子识别Markov毯并说明消去项？
  \item 能否为原书五分类任务补齐先验、推导精确条件混合并用解析矩核验运行？若不能，重做练习\ref{ex:V5-C04-S04-01}。
  \item 能否拒绝用接受率、燃烧期或单链图替代不可约性和程序正确性证据？若不能，重做练习\ref{ex:V5-C04-S04-04}。
\end{enumerate}
\end{selfcheckbox}
\end{SLCompactClosure}
